\documentclass[11pt,a4paper]{article}

\usepackage[utf8]{inputenc}
\usepackage[T1]{fontenc}
\usepackage[turkish]{babel}
\usepackage[margin=2.3cm]{geometry}
\usepackage{titlesec}
\usepackage{enumitem}
\usepackage{xcolor}
\usepackage{hyperref}
\usepackage{fancyhdr}
\usepackage{parskip}
\usepackage{listings}

\definecolor{basliknoktu}{RGB}{20,60,110}
\definecolor{kutucerceve}{RGB}{60,60,60}
\definecolor{kutuarkaplan}{RGB}{245,245,245}
\definecolor{kodarkaplan}{RGB}{40,42,54}
\definecolor{kodyazi}{RGB}{248,248,242}
\definecolor{kodyorum}{RGB}{150,150,160}
\definecolor{kodanahtar}{RGB}{255,184,108}
\definecolor{linkrengi}{RGB}{20,90,140}

\titleformat{\section}
  {\normalfont\Large\bfseries\color{basliknoktu}}
  {B\"ol\"um \thesection}{0.6em}{}
\titleformat{\subsection}
  {\normalfont\large\bfseries}
  {}{0em}{}
\titleformat{\subsubsection}
  {\normalfont\normalsize\bfseries\itshape}
  {}{0em}{}

\hypersetup{
  colorlinks=true,
  linkcolor=linkrengi,
  urlcolor=linkrengi,
  citecolor=linkrengi,
  pdftitle={LDD3 Calisma Kagidi - Detayli},
  pdfauthor={Ahmet Hasan Celik}
}

\pagestyle{fancy}
\fancyhf{}
\fancyhead[L]{LDD3 -- Detayl\i\ \c{C}al\i\c{s}ma Ka\u{g}\i d\i\ (B\"ol\"um 2-6)}
\fancyhead[R]{\thepage}
\renewcommand{\headrulewidth}{0.4pt}

\lstdefinestyle{kernelc}{
  backgroundcolor=\color{kodarkaplan},
  basicstyle=\ttfamily\footnotesize\color{kodyazi},
  commentstyle=\color{kodyorum}\itshape,
  keywordstyle=\color{kodanahtar}\bfseries,
  numberstyle=\tiny\color{kodyorum},
  breaklines=true,
  breakatwhitespace=false,
  frame=single,
  rulecolor=\color{kutucerceve},
  framesep=4pt,
  xleftmargin=4pt,
  xrightmargin=4pt,
  showstringspaces=false,
  tabsize=4,
  morekeywords={struct,static,const,unsigned,typedef,int,void,char,long,size_t,ssize_t,loff_t,dev_t}
}
\lstset{style=kernelc}

\newenvironment{notkutu}[1]{%
  \par\medskip\noindent
  \begin{minipage}{\linewidth}
  \begin{center}
  \fcolorbox{kutucerceve}{kutuarkaplan}{%
    \begin{minipage}{0.95\linewidth}
    \vspace{2pt}
    \textbf{#1}\par\smallskip
}}{%
    \vspace{2pt}
    \end{minipage}%
  }
  \end{center}
  \end{minipage}\par\medskip
}

\newenvironment{ahmetkutu}{%
  \par\medskip\noindent
  \begin{minipage}{\linewidth}
  \begin{center}
  \fcolorbox{basliknoktu}{kutuarkaplan}{%
    \begin{minipage}{0.95\linewidth}
    \vspace{2pt}
    \textbf{ahmet.c ile ba\u{g}lant\i}\par\smallskip
}}{%
    \vspace{2pt}
    \end{minipage}%
  }
  \end{center}
  \end{minipage}\par\medskip
}

\title{\Huge Linux Device Drivers (3. Bask\i)\\[0.3em]
\Large Detayl\i\ \c{C}al\i\c{s}ma Ka\u{g}\i d\i\ -- B\"ol\"um 2--6, 8, 10, 14}
\author{Ahmet Hasan \c{C}elik \\ \small G\"om\"ulu Sistem Staj\i\ -- Kernel Modu\"l\"u Geli\c{s}tirme}
\date{4 A\u{g}ustos 2026}

\begin{document}

\maketitle

\begin{notkutu}{Bu belge nedir?}
Bu, LDD3'ün ilk özet çalışma kağıdının \emph{derinleştirilmiş} sürümüdür.
İlk sürümdeki 18 bölümlük kısa özet yerine, burada staj için en kritik
sekiz bölüm -- \textbf{2 (Modülleri Derleme/Çalıştırma), 3 (Karakter
Sürücüler), 4 (Hata Ayıklama), 5 (Eşzamanlılık), 6 (Gelişmiş Karakter
Sürücü İşlemleri), 8 (Bellek Ayırma), 10 (Kesme Yönetimi), 14 (Linux
Aygıt Modeli)} -- ele alınıyor, ama her biri kitabın kendisinden
gerçekten okunarak, fonksiyon imzalarıyla, gerçek \texttt{scull} sürücüsü
koduyla ve kitabın açıkça belirttiği tuzaklarla (gotcha) birlikte
anlatılıyor.

Amaç: bu belgeyi okuduktan sonra kitabın ilgili sayfalarını açtığında
hiçbir kavramın "havada kalmaması" ve kendi \texttt{ahmet.c} sürücünle
aradaki farkı/benzerliği net görebilmen.
\end{notkutu}

\tableofcontents

\newpage

%=================================================================
\section{Mod\"ulleri Derleme ve \c{C}al\i\c{s}t\i rma}
\label{sec:bolum2}
%=================================================================

\subsection{Hello World: Bir Mod\"ul\"un Anatomisi}

Kitap ilk örneğini şu şekilde verir:

\begin{lstlisting}
#include <linux/init.h>
#include <linux/module.h>
MODULE_LICENSE("Dual BSD/GPL");

static int hello_init(void)
{
    printk(KERN_ALERT "Hello, world\n");
    return 0;
}

static void hello_exit(void)
{
    printk(KERN_ALERT "Goodbye, cruel world\n");
}

module_init(hello_init);
module_exit(hello_exit);
\end{lstlisting}

\textbf{Kritik nokta:} fonksiyonların adı \texttt{hello\_init}/\texttt{hello\_exit}
olması bir zorunluluk değil -- \texttt{module\_init()} ve \texttt{module\_exit()}
birer \emph{makro}dır ve kernel'e "yüklenince/kaldırılınca hangi
fonksiyonu çağıracağını" söyler. Sen bu fonksiyonlara istediğin ismi
verebilirsin (senin kodunda \texttt{ahmet\_init}/\texttt{ahmet\_exit}).

\begin{notkutu}{Beklenmedik ama önemli bir ayrıntı}
Kitap açıkça bir yazım hatası tuzağına dikkat çekiyor: \texttt{KERN\_ALERT}
ile format string'i arasında \textbf{virgül yoktur}. \texttt{printk(KERN\_ALERT,
"...")} yazmak yaygın bir acemi hatasıdır -- \texttt{KERN\_ALERT} aslında
bir fonksiyon parametresi değil, derleme zamanında format string'in
başına \emph{yapıştırılan} bir string sabittir (\texttt{"<1>"} gibi).
Yani doğrusu \texttt{printk(KERN\_ALERT "mesaj\textbackslash n")} --
iki string derleyici tarafından otomatik birleştirilir (string
concatenation). Neyse ki derleyici virgülü koyarsan hata verir, bu
yüzden tehlike sınırlı.
\end{notkutu}

\subsection{Neden \texttt{printf} değil de \texttt{printk}?}

Kernel'de standart C kütüphanesi (glibc) yoktur -- modülün bağlı olduğu
tek şey kernel'in kendisidir. \texttt{printk} kernelin kendi
\texttt{printf} benzeri fonksiyonudur (kayan nokta -- floating point --
desteği \textbf{yoktur}, çünkü kernel modunda FP register durumunu
kaydetme/geri yükleme maliyeti gereksiz bulunmuştur). \texttt{printk}'in
\texttt{printf}'ten en büyük farkı: bir \textbf{öncelik seviyesi}
(loglevel) alabilmesidir -- bunu Bölüm 4'te detaylı göreceğiz.

\subsection{Uygulama ile Mod\"ul Aras\i ndaki Fark: Olay G\"ud\"uml\"u Model}

Bir C uygulaması \texttt{main()}'den başlar, sırayla çalışır, biter. Bir
kernel modülünün ise \texttt{main()}'i yoktur. \texttt{init} fonksiyonu
"ben buradayım, şunları yapabilirim" der ve kaynakları/yeteneklerini
\textbf{kaydedip hemen döner} -- modülün asıl işi, kayıt ettiği
fonksiyonlar (senin \texttt{fops} tablon gibi) başka biri tarafından
\emph{daha sonra} çağrıldığında yapılır. \texttt{exit} fonksiyonu ise
"artık burada değilim" der, kayıt olan her şeyi geri alır.

\begin{notkutu}{Önemli fark: temizlik zorunluluğu}
Bir uygulama çöktüğünde/bittiğinde işletim sistemi onun kaynaklarını
otomatik temizler (dosya tanıtıcıları, bellek...). Bir kernel modülünün
\texttt{exit} fonksiyonu ise \textbf{kendi açtığı her şeyi kendisi elle
kapatmak zorundadır} -- kernel, "bu modül şunları kaydetti" diye bir
kayıt tutmaz. Bir kaynağı serbest bırakmayı unutursan, o kaynak sistem
yeniden başlatılana kadar (ya da bir sonraki modül yüklemesi
çakışana kadar) sızdırılmış kalır.
\end{notkutu}

\subsection{Kernel Stack'inin K\"u\c{c}\"ukl\"u\u{g}\"u}

Kitabın açıkça vurguladığı bir tuzak: kernel stack'i çok küçüktür --
bazı mimarilerde tek bir 4096 byte'lık sayfa kadar, ve bu stack
\textbf{tüm kernel-alanı çağrı zincirinde paylaşılır}. Bu yüzden kernel
kodunda \textbf{büyük yerel (otomatik/stack) değişkenler asla
tanımlanmamalıdır} -- örneğin fonksiyon içinde \texttt{char buf[8192];}
gibi bir tanım kernel'de ciddi bir hatadır (stack overflow riski);
büyük veri için \texttt{kmalloc} (Bölüm 8) kullanılmalıdır.

\subsection{current: \c{S}u Anki Process}

\texttt{current} -- \texttt{<asm/current.h>} üzerinden erişilen, o an
çalışan process'i temsil eden \texttt{struct task\_struct *} tipinde bir
global işaretçi (aslında CPU başına ayrı tutulan, mimariye özel bir
mekanizma, düz bir global değişken değil). Örnek:

\begin{lstlisting}
printk(KERN_INFO "The process is \"%s\" (pid %i)\n",
       current->comm, current->pid);
\end{lstlisting}

\texttt{current->comm} -- çalıştıran programın adı (15 karaktere
kırpılmış); \texttt{current->pid} -- process ID. Bunu Bölüm 6'da
\texttt{sculluid} örneğinde (kullanıcı bazlı erişim kontrolü) tekrar
göreceksin.

\subsection{Kbuild: Kernel Modülü Neden Normal \texttt{gcc} ile Derlenmez}

Modülün, derlendiği kernel ile \textbf{birebir aynı ABI'ye, aynı yapı
hizalamalarına} sahip olması gerekir -- bu yüzden kendi başına
\texttt{gcc dosya.c} ile derlenemez, kernel'in kendi build sistemi
(\emph{Kbuild}) kullanılır.

En basit \texttt{obj-m} kullanımı:
\begin{lstlisting}
obj-m := hello.o
\end{lstlisting}

Birden fazla kaynak dosyadan tek modül üretmek istersen:
\begin{lstlisting}
obj-m := module.o
module-objs := file1.o file2.o
\end{lstlisting}

Gerçek çağrı, kernel ağacının makefile'ını kullanarak yapılır:
\begin{lstlisting}
make -C ~/kernel-2.6 M=`pwd` modules
\end{lstlisting}
\texttt{-C} kernel kaynak dizinine geçer (üst-düzey makefile'ı bulur);
\texttt{M=} o makefile'a "işini bitirince modül dizinime geri dön ve
\texttt{obj-m} listesini kullanarak \texttt{modules} hedefini derle"
der.

Kitabın önerdiği, kendini iki kez okuyan (bir kere doğrudan, bir kere
Kbuild tarafından \texttt{KERNELRELEASE} tanımlıyken) idiyomatik
makefile kalıbı:
\begin{lstlisting}
ifneq ($(KERNELRELEASE),)
    obj-m := hello.o
else
    KERNELDIR ?= /lib/modules/$(shell uname -r)/build
    PWD := $(shell pwd)
default:
    $(MAKE) -C $(KERNELDIR) M=$(PWD) modules
endif
\end{lstlisting}

\begin{ahmetkutu}
Senin \texttt{Makefile}'ın bu kalıbın basitleştirilmiş hali:
\begin{lstlisting}
obj-m += ahmet.o
KDIR := /lib/modules/$(shell uname -r)/build
PWD := $(shell pwd)
all:
	$(MAKE) -C $(KDIR) M=$(PWD) modules
clean:
	$(MAKE) -C $(KDIR) M=$(PWD) clean
\end{lstlisting}
\texttt{KERNELRELEASE} kontrolü olmadan, çünkü sen hep doğrudan
komut satırından \texttt{make} çağırıyorsun -- ama sonuç birebir aynı
mantık: derleme işi kernel ağacındaki gerçek build sistemine devredilir.
\end{ahmetkutu}

\subsection{insmod, modprobe, rmmod, lsmod -- Ger\c{c}ekte Ne Yaparlar}

\begin{itemize}[leftmargin=1.2cm]
  \item \textbf{insmod} -- modül kodunu+verisini kernel belleğine
  yükler; bir \emph{linker} gibi davranır, modüldeki çözümlenmemiş
  sembolleri kernel'in sembol tablosuna karşı bağlar. Disk üzerindeki
  dosyayı değiştirmez, bellekteki bir kopya üzerinde çalışır. İçeride
  \texttt{sys\_init\_module} sistem çağrısını tetikler: kernel
  \texttt{vmalloc} ile bellek ayırır (Bölüm 8), modül metnini oraya
  kopyalar, kernel sembol tablosuna karşı referansları çözer, ve
  modülün init fonksiyonunu çağırır.
  \item \textbf{modprobe} -- \texttt{insmod} gibi yükler, ama ek olarak
  modülün ihtiyaç duyduğu \textbf{başka modülleri de otomatik bulup
  yükler} (bağımlılık çözümü). Salt \texttt{insmod} kullanırsan ve
  eksik sembol varsa, sistem log dosyasında "unresolved symbols"
  hatası alırsın.
  \item \textbf{rmmod} -- modülü kaldırır; modül hâlâ kullanımda
  görünüyorsa (örn. açık bir dosya tanıtıcısı ona referans veriyorsa)
  başarısız olur. Kitap açıkça uyarıyor: "zorla kaldırma" seçeneği
  olsa da, o noktaya gelmişsen muhtemelen işler zaten yeterince
  bozulmuştur -- \textbf{yeniden başlatmak daha güvenli bir yoldur.}
  \item \textbf{lsmod} -- yüklü modülleri listeler; \texttt{/proc/modules}
  dosyasını okuyarak çalışır (aynı bilgi \texttt{/sys/module} altında
  da mevcuttur).
\end{itemize}

\subsection{Modül Parametreleri}

Bir sürücünün sisteme özel bilgiye ihtiyacı olabilir (I/O adresi, DMA
kanalı, bir tampon boyutu...) -- bunları derleme zamanında sabitlemek
yerine \textbf{yükleme zamanında} (\texttt{insmod}/\texttt{modprobe}
ile) ayarlanabilir hale getirmek için \texttt{module\_param} makrosu
kullanılır:

\begin{lstlisting}
#include <linux/moduleparam.h>

static char *whom = "world";
static int howmany = 1;
module_param(howmany, int, S_IRUGO);
module_param(whom, charp, S_IRUGO);
\end{lstlisting}

\begin{lstlisting}
insmod hellop howmany=10 whom="Mom"
\end{lstlisting}

\begin{itemize}[leftmargin=1.2cm]
  \item \textbf{Desteklenen tipler}: \texttt{bool}, \texttt{invbool}
  (tersine çevrilmiş boolean), \texttt{charp} (char pointer -- string),
  \texttt{int}, \texttt{long}, \texttt{short}, ve işaretsiz
  (\texttt{u}-önekli) karşılıkları.
  \item \textbf{Dizi parametreleri}: \texttt{module\_param\_array(name,
  type, num, perm)} -- \texttt{num}, kaç değer girildiğini kaydeden bir
  \texttt{int} işaretçisi.
  \item \textbf{\texttt{perm} (izin) argümanı} -- \texttt{<linux/stat.h>}
  içindeki tanımları kullanır, parametrenin \texttt{/sys/module/<isim>/
  parameters/} altındaki dosyasının izinlerini belirler:
  \begin{itemize}
    \item \texttt{0} $\to$ sysfs'te hiç görünmez.
    \item \texttt{S\_IRUGO} $\to$ herkes okuyabilir, değiştirilemez.
    \item \texttt{S\_IRUGO|S\_IWUSR} $\to$ root sysfs üzerinden
    değiştirebilir.
  \end{itemize}
\end{itemize}

\begin{notkutu}{Gizli tuzak: yazılabilir parametreler}
Bir parametre sysfs üzerinden değiştirilebilir yapılırsa, modülün
gördüğü değer gerçekten değişir -- \textbf{ama modül bu değişiklikten
başka hiçbir şekilde haberdar edilmez} (bir callback/interrupt
tetiklenmez). Bu yüzden bir parametreyi yazılabilir yapmadan önce,
değişikliği algılayıp tepki verecek bir mekanizman olduğundan emin
olmalısın -- yoksa parametreyi yazılabilir yapmamalısın.
\end{notkutu}

\begin{ahmetkutu}
Senin \texttt{ahmet.c}'nde \texttt{\#define BUFFER\_SIZE 1024} sabit
bir derleme-zamanı değeri. Bunu \texttt{module\_param(buffer\_size, int,
S\_IRUGO)} ile yükleme-zamanı parametresine çevirmek, bu bölümün
öğrettiği en doğrudan uygulanabilir tekniktir -- README'deki "sırada"
listesine bu netlikle eklenebilir.
\end{ahmetkutu}

\subsection{Kernel Sembol Tablosu ve Mod\"ul Y\i\u{g}\i lmas\i\ (Stacking)}

Bir modül başka bir modülün sembollerini kullanabilir -- buna
\textbf{modül yığılması} denir. Gerçek örnekler: \texttt{msdos}
dosya sistemi \texttt{fat} modülünün sembollerini kullanır; paralel
port alt sistemi \texttt{lp} $\to$ \texttt{parport} $\to$
\texttt{parport\_pc} şeklinde katmanlanır. Bir sembolü dışa açmak için:

\begin{lstlisting}
EXPORT_SYMBOL(isim);
EXPORT_SYMBOL_GPL(isim);
\end{lstlisting}

İkincisi sembolü sadece GPL lisanslı modüllere açar. Çoğu modülün hiç
sembol dışa açmasına gerek yoktur -- sadece başka modüllerin
gerçekten faydalanacağı durumlarda kullanılır.

\subsection{Hata Y\"onetimi: goto Kal\i b\i\ ve S\"okme S\i ras\i}

Kayıt (registration) işlemleri başarısız olabilir (bellek yetersizliği
gibi). Linux, hangi modülün neyi kaydettiğine dair \textbf{merkezi bir
defter tutmaz} -- yani bir kayıt başarısız olduğunda, o ana kadar
başarıyla yapılmış olan \emph{önceki} kayıtları geri almak (unregister)
tamamen senin sorumluluğundadır. Kitap, iç içe \texttt{if}'lerden
kaçınmak için normalde sevilmeyen \texttt{goto}'yu burada
\textbf{bilinçli olarak önerir}:

\begin{lstlisting}
int __init my_init_function(void)
{
    int err;
    err = register_this(ptr1, "skull");
    if (err) goto fail_this;
    err = register_that(ptr2, "skull");
    if (err) goto fail_that;
    err = register_those(ptr3, "skull");
    if (err) goto fail_those;
    return 0; /* basari */

  fail_those: unregister_that(ptr2, "skull");
  fail_that:  unregister_this(ptr1, "skull");
  fail_this:  return err;
}
\end{lstlisting}

Cleanup fonksiyonu, kayıtları \textbf{ters sırada} söker (A, B, C
kaydedildiyse; C, B, A sökülür) -- bu, kernel programlamada genel bir
prensiptir ve senin \texttt{ahmet\_exit}'inde de
(\texttt{device\_destroy} $\to$ \texttt{class\_destroy} $\to$
\texttt{cdev\_del} $\to$ \texttt{unregister\_chrdev\_region}) aynen
uygulanmış durumda.

\begin{notkutu}{İki önemli race-condition uyarısı (detayları Bölüm 5'te)}
\begin{enumerate}[leftmargin=0.8cm]
  \item Bir modül bir yeteneği kaydettiği \textbf{anda}, kernelin başka
  bir parçası onu hemen kullanmaya başlayabilir -- init fonksiyonun
  hâlâ çalışıyor olması bunu engellemez. \textbf{Kural: bir yeteneği,
  onu desteklemek için gereken tüm iç hazırlık tamamlanmadan
  kaydetme.}
  \item Eğer init, bir yetenek zaten kaydedilip \emph{kullanılmaya
  başlandıktan sonra} başarısız olursa, bu senaryoyu düşün: belki de
  hiç başarısız olmamak (init'i bitirmek) daha güvenlidir.
\end{enumerate}
\end{notkutu}

\subsection{Kullanıcı Alanında mı, Kernel'de mi Yazmalı? (Doing It in User Space)}

Kitap, çekingen kernel programcılarına şu soruyu sorduruyor: her zaman
kernel modülü yazmak mı gerekir?

\textbf{Kullanıcı alanı sürücüsünün avantajları} (kitaptan, aynen):
\begin{enumerate}[leftmargin=1.2cm]
  \item Tam C kütüphanesi kullanılabilir.
  \item Normal bir debugger (gdb) rahatça kullanılabilir.
  \item Sürücü kilitlenirse sadece o process öldürülür, sistem
  genelde etkilenmez.
  \item Kullanıcı belleği \emph{swap edilebilir} -- kernel belleği
  edilemez; nadiren kullanılan büyük bir sürücü sürekli RAM işgal
  etmez.
  \item Kapalı kaynaklı sürücüler için lisans belirsizliği ortadan
  kalkar.
\end{enumerate}

\textbf{Dezavantajları}:
\begin{enumerate}[leftmargin=1.2cm]
  \item Kesmeler (interrupt) user space'te mevcut değildir.
  \item Doğrudan bellek erişimi sadece \texttt{/dev/mem}'i
  \texttt{mmap} ederek mümkün, ve bunu sadece yetkili kullanıcı
  yapabilir.
  \item I/O port erişimi kısıtlı ve/veya yavaş.
  \item Yanıt süresi daha yavaştır (context switch gerekir); sürücü
  disk'e swap edilmişse çok daha kötüleşir.
  \item \textbf{En önemli cihaz sınıfları user space'te hiç
  yapılamaz}: ağ arayüzleri ve blok cihazlar bunlara dahildir.
\end{enumerate}

Kitabın tavsiyesi: yeni/alışılmadık bir donanımı ilk öğrenirken user
space'te denemek, tüm sistemi riske atmadan deney yapmayı sağlar --
donanım anlaşıldıktan sonra "yazılımı bir kernel modülüne kapsüllemek
sıkıntısız bir işlem olmalı."

\subsection*{B\"ol\"um 2 -- H\i zl\i\ Referans (kitab\i n kendi \"ozeti)}
\begin{itemize}[leftmargin=1.2cm]
  \item \texttt{insmod}, \texttt{modprobe}, \texttt{rmmod} -- modül
  yükleme/kaldırma araçları.
  \item \texttt{module\_init(f)}, \texttt{module\_exit(f)} -- init/exit
  fonksiyonlarını bildiren makrolar.
  \item \texttt{\_\_init}, \texttt{\_\_initdata}, \texttt{\_\_exit},
  \texttt{\_\_exitdata} -- yalnızca init/exit zamanında kullanılan
  fonksiyon/veri işaretleri; bellekten atılabilirler (özel bir ELF
  bölümüne yerleştirilerek).
  \item \texttt{struct task\_struct *current;} -- şu anki process.
  \item \texttt{obj-m} -- Kbuild'e hangi modüllerin derleneceğini
  söyleyen makefile sembolü.
  \item \texttt{/sys/module}, \texttt{/proc/modules} -- yüklü modüller
  hakkında bilgi.
  \item \texttt{EXPORT\_SYMBOL(sym)}, \texttt{EXPORT\_SYMBOL\_GPL(sym)}
  -- bir sembolü kernel'e (veya sadece GPL modüllere) açar.
  \item \texttt{module\_param(degisken, tip, izin)} -- yükleme-zamanı
  ayarlanabilir bir parametre tanımlar.
  \item \texttt{int printk(const char *fmt, ...);} -- kernelin
  \texttt{printf} karşılığı.
\end{itemize}

\newpage

%=================================================================
\section{Karakter S\"ur\"uc\"uler}
\label{sec:bolum3}
%=================================================================

Bu bölüm kitabın en yoğun kod içeren bölümlerinden biridir çünkü artık
gerçek bir sürücü (\texttt{scull}) baştan sona inşa edilir. Senin
\texttt{ahmet.c}'n bu bölümün kavramlarının basitleştirilmiş bir
uygulamasıdır -- burada scull'un tam mimarisini görüp kendi kodunla
karşılaştıracaksın.

\subsection{scull'un Tasar\i m\i : Neden Birden Fazla ``Ki\c{s}ilik''?}

\texttt{scull} = \emph{Simple Character Utility for Loading Localities}.
Gerçek donanıma değil, kernel belleğinden ayrılmış bir alana çalışır --
bu yüzden taşınabilir, herkes derleyip çalıştırabilir. Kitap tek bir
mekanizma üzerine \textbf{altı farklı ``kişilik''} inşa eder, her biri
gerçek dünyadaki bir erişim politikasını örnekler:

\begin{itemize}[leftmargin=1.2cm]
  \item \textbf{scull0--scull3}: dört cihaz, her biri hem
  \emph{global} (aynı dosyayı açan tüm süreçler aynı veriyi paylaşır --
  senin \texttt{ahmet}'inin de yaptığı gibi) hem \emph{kalıcı} (veri
  kapat/aç arasında kaybolmaz, sadece açık truncate ile silinir).
  \item \textbf{scullpipe0--3}: FIFO/boru gibi davranan dört cihaz --
  Bölüm 6'da bloklayan I/O örneği için kullanılır.
  \item \textbf{scullsingle}: scull0 gibi ama \textbf{aynı anda sadece
  bir process} açabilir.
  \item \textbf{scullpriv}: her sanal konsol/terminale \textbf{özel}
  bir bellek alanı.
  \item \textbf{sculluid}: aynı anda sadece \textbf{bir kullanıcı}
  (UID) açabilir; ikinci kullanıcı hata alır.
  \item \textbf{scullwuid}: sculluid gibi ama hata vermek yerine
  \textbf{bekler} (bloklar).
\end{itemize}

Kitabın kendi ifadesiyle: bu erişim-kontrollü varyantlar "mekanizma ile
politikayı karıştırıyormuş gibi görünebilir, ama gerçek hayattaki bazı
cihazlar tam olarak bu tür yönetim gerektirir." Bu bölüm sadece
scull0--3'ün iç yapısını kapsar; scullpipe ve erişim-kontrollü
varyantlar Bölüm 6'da.

\subsection{Major/Minor Numaralar\i n\i n \.{I}\c{c} Temsili}

\texttt{dev\_t} tipi (\texttt{<linux/types.h>}), major ve minor'u tek
bir sayıda taşır. 2.6.0'dan itibaren: 32 bit, \textbf{12 bit major, 20
bit minor}. Bu düzene asla doğrudan (bit kaydırma ile) güvenme --
her zaman şu makroları kullan (\texttt{<linux/kdev\_t.h>}):

\begin{lstlisting}
MAJOR(dev_t dev);        /* major'u cikar */
MINOR(dev_t dev);        /* minor'u cikar */
MKDEV(int major, int minor);  /* ikisinden dev_t olustur */
\end{lstlisting}

Kitap açıkça uyarıyor: \texttt{dev\_t}'nin iç yapısı gelecekte tekrar
değişebilir -- taşınabilir kod bu makrolarla yazılmalı, ham bit
işlemleriyle değil.

\subsection{Cihaz Numaras\i\ Ay\i rma: \texttt{register\_chrdev\_region} vs \texttt{alloc\_chrdev\_region}}

\begin{lstlisting}
int register_chrdev_region(dev_t first, unsigned int count, char *name);
int alloc_chrdev_region(dev_t *dev, unsigned int firstminor,
                         unsigned int count, char *name);
void unregister_chrdev_region(dev_t first, unsigned int count);
\end{lstlisting}

\begin{itemize}[leftmargin=1.2cm]
  \item \texttt{register\_chrdev\_region}: major numarasını \textbf{sen
  zaten biliyorsan} kullanılır (\texttt{first} parametresiyle).
  \item \texttt{alloc\_chrdev\_region}: \texttt{dev} bir \textbf{çıktı}
  parametresidir -- kernel dinamik olarak bir major seçer ve sonucu
  oraya yazar. Kitap açıkça şöyle diyor: \textbf{"yeni sürücülerin
  neredeyse kesinlikle \texttt{alloc\_chrdev\_region} kullanması
  gerekir"} -- rastgele boş bir major seçmek "sadece sürücünün tek
  kullanıcısı sen olduğun sürece işe yarar," sürücü yaygınlaştıkça
  çakışma riski artar.
\end{itemize}

\begin{notkutu}{Dinamik atamanın bedeli}
Major dinamik olduğu için \texttt{/dev} düğümlerini önceden
oluşturamazsın -- yükledikten sonra \texttt{/proc/devices}'ten atanan
major'u okuyup (genelde bir yükleme betiğinde \texttt{awk} ile) ona
göre \texttt{mknod} yapman gerekir. \textbf{Ama senin \texttt{ahmet.c}'in
bu problemi tamamen ortadan kaldırıyor} -- çünkü Bölüm 14'teki Linux
Device Model'i (\texttt{class\_create}/\texttt{device\_create})
kullanıyorsun; udev major'u otomatik öğrenip \texttt{/dev/ahmet}'i
kendisi oluşturuyor. Kitabın scull'u bu kolaylığı henüz kullanmıyor
(scull, Bölüm 3'te yazıldığı haliyle, elle \texttt{mknod} betikleriyle
çalışır) -- senin sürücün bu noktada scull'dan daha modern.
\end{notkutu}

\subsection{\"U\c{c} Temel Veri Yap\i s\i: file\_operations, file, inode}

\subsubsection{struct file\_operations}

Cihaz numaralarını sürücü koduna bağlayan tablo. Her açık dosyanın
\texttt{f\_op} alanı kendi \texttt{file\_operations}'ına işaret eder.
\texttt{\_\_user} etiketi (örn. \texttt{char \_\_user *buf}) bir
işaretçinin \textbf{user-space adresi} olduğunu, kernel kodunun onu
doğrudan dereference edemeyeceğini belgeler (statik analiz araçları
için, derleyici için etkisi yoktur).

En çok kullanılan alanlar (NULL bırakılırsa ne olduğuyla birlikte):

\begin{itemize}[leftmargin=1.2cm]
  \item \texttt{owner} -- işlem değil, \texttt{THIS\_MODULE} olmalı;
  kullanımdayken modülün kaldırılmasını engeller.
  \item \texttt{read}/\texttt{write} -- NULL ise syscall
  \texttt{-EINVAL} döner.
  \item \texttt{ioctl} -- NULL ise öntanımlı olmayan her ioctl
  \texttt{-ENOTTY} döner.
  \item \texttt{open} -- NULL ise açma her zaman başarılı olur ama
  sürücü haberdar edilmez.
  \item \texttt{release} -- NULL olabilir; \textbf{her \texttt{close()}
  çağrısında değil}, sadece \texttt{file} yapısı gerçekten yok
  edildiğinde çağrılır (aşağıda detaylı).
  \item \texttt{llseek} -- NULL ise kernel \texttt{f\_pos}'u kendi
  varsayılan mantığıyla değiştirir (Bölüm 6'da detaylı).
  \item \texttt{mmap}, \texttt{poll}, \texttt{fasync}, \texttt{fsync}
  -- ileri düzey; sırasıyla Bölüm 15 ve Bölüm 6'da.
\end{itemize}

scull'un gerçek tanımı, C99 adlandırılmış alan (designated initializer)
söz dizimiyle:
\begin{lstlisting}
struct file_operations scull_fops = {
    .owner =   THIS_MODULE,
    .llseek =  scull_llseek,
    .read =    scull_read,
    .write =   scull_write,
    .ioctl =   scull_ioctl,
    .open =    scull_open,
    .release = scull_release,
};
\end{lstlisting}
Bu söz dizimi tercih edilir çünkü: (1) struct'ın alan sırası değişse
bile kod kırılmaz, (2) belirtilmeyen alanlar otomatik NULL olur, (3)
sık erişilen işaretçiler aynı cache satırına yerleştirilebilir.

\subsubsection{struct file: f\_pos, f\_flags, private\_data}

\texttt{struct file} \textbf{açık bir dosyayı} temsil eder -- C
kütüphanesindeki \texttt{FILE*} ile \textbf{karıştırılmamalı}, ikisi
birbirinden tamamen bağımsızdır. Sürücü yazarları için önemli alanlar:

\begin{itemize}[leftmargin=1.2cm]
  \item \texttt{f\_pos} (\texttt{loff\_t}, 64-bit) -- güncel okuma/yazma
  konumu. Sürücü bunu \textbf{okuyabilir} ama normalde
  \textbf{doğrudan yazmamalı} -- konum güncellemesi \texttt{read}/
  \texttt{write}'a geçirilen \texttt{*offp} işaretçisi üzerinden
  yapılmalı (istisna: \texttt{llseek}'in tüm işi zaten konum
  değiştirmektir).
  \item \texttt{f\_flags} -- \texttt{O\_RDONLY}, \texttt{O\_NONBLOCK},
  \texttt{O\_SYNC} gibi bayraklar. \textbf{Okuma/yazma izni kontrolü
  için \texttt{f\_mode} kullanılmalı, \texttt{f\_flags} değil.}
  \item \texttt{f\_op} -- kernel tarafından \texttt{open} sırasında
  atanır. \textbf{Kernel bu değeri başka hiçbir yerde önbelleğe
  almaz} -- yani bir sürücü kendi \texttt{open}'ı içinde
  \texttt{filp->f\_op}'u \textbf{farklı bir \texttt{file\_operations}
  tablosuna değiştirebilir}, ve bu değişiklik o dosya tanıtıcısındaki
  sonraki tüm çağrılarda geçerli olur ("method overriding" -- kernelin
  nesne-yönelimli hilesi). \texttt{/dev/null} ve \texttt{/dev/zero}
  (ikisi de major 1) aynı major altında farklı davranışlar sunmak için
  bunu kullanır.
  \item \texttt{private\_data} -- \texttt{open} syscall'ı tarafından
  \texttt{NULL}'a ayarlanır, sürücü serbestçe kullanabilir. \textbf{Bu
  senin \texttt{ahmet.c}'nin şu an kullanmadığı ama scull'un aktif
  kullandığı alan} -- scull, \texttt{open} içinde kendi
  \texttt{scull\_dev} yapısına işaretçiyi buraya koyar, böylece
  \texttt{read}/\texttt{write} her seferinde hangi cihaz örneğiyle
  uğraştığını bilir.
\end{itemize}

\begin{notkutu}{release vs. flush: gizli bir tuzak}
\texttt{release}, her \texttt{close()} çağrısında \textbf{çağrılmaz}.
\texttt{fork()} ve \texttt{dup()}, \emph{aynı} \texttt{file} yapısına
işaret eden yeni dosya tanıtıcı \textbf{kopyaları} oluşturur --
\texttt{open()} çağırmazlar. Kernel, \texttt{file} yapısı üzerinde bir
\textbf{kullanım sayacı} tutar; \texttt{close()} sadece bu sayaç
\textbf{sıfıra düştüğünde} \texttt{release}'i çağırır -- yani her
\texttt{open} için tam olarak bir \texttt{release} garanti edilir.
Buna karşılık \texttt{flush}, uygulama her \texttt{close} çağırdığında
(her fd kopyası için) tetiklenir -- ama çoğu sürücü bunu
uygulamaz çünkü genelde \texttt{close} anında (asıl \texttt{release}
olmadıkça) yapılacak bir şey yoktur.
\end{notkutu}

\subsection{cdev: Sürücüyü Kernel'e Tanıtmak}

\begin{lstlisting}
void cdev_init(struct cdev *cdev, struct file_operations *fops);
int  cdev_add(struct cdev *dev, dev_t num, unsigned int count);
void cdev_del(struct cdev *dev);
\end{lstlisting}

İki kullanım yolu: (1) \texttt{cdev\_alloc()} ile bağımsız tahsis, ya
da (2) \textbf{scull'un yaptığı gibi} \texttt{cdev}'i kendi cihaz
struct'ının \textbf{içine gömmek}:

\begin{lstlisting}
struct scull_dev {
    struct scull_qset *data;  /* ilk quantum set'e isaretci */
    int quantum;              /* guncel quantum boyutu */
    int qset;                 /* guncel dizi boyutu */
    unsigned long size;       /* burada saklanan veri miktari */
    unsigned int access_key;  /* sculluid/scullpriv icin */
    struct semaphore sem;     /* karsilikli dislama semaforu */
    struct cdev cdev;         /* Char device yapisi */
};
\end{lstlisting}

\begin{ahmetkutu}
Senin kodunda \texttt{static struct cdev ahmet\_cdev;} ayrı, bağımsız
bir global değişken -- scull'un gömme (embedding) desenini
\textbf{kullanmıyorsun}. Bu, tek cihazlı basit bir sürücü için
tamamen makul bir sadeleştirme. Eğer ileride birden fazla \texttt{ahmet}
cihazı (minor 0, 1, 2...) desteklemek istersen, scull'un yaptığı gibi
\texttt{cdev}'i kendi struct'ının içine gömüp bir dizi/liste
oluşturman gerekecek -- tam olarak bu bölümün öğrettiği desen.
\end{ahmetkutu}

\textbf{Kritik sıralama kuralı}: \texttt{cdev\_add()} döndüğü anda
cihaz "canlıdır" -- kernel operasyonlarını \textbf{hemen} çağırabilir.
Bu yüzden \texttt{cdev\_add}, sürücü tamamen hazır olmadan
\textbf{asla} çağrılmamalı.

\subsubsection{container\_of: inode'dan kendi struct'\i na geri d\"onmek}

scull'un \texttt{open}'ı, \texttt{inode->i\_cdev} üzerinden embed
edilmiş \texttt{cdev}'e erişebilir -- ama asıl istenen \emph{onu içeren}
\texttt{scull\_dev}'dir. Çözüm, \texttt{<linux/kernel.h>}'deki genel
amaçlı bir makro:

\begin{lstlisting}
container_of(pointer, container_type, container_field);
\end{lstlisting}

\begin{lstlisting}
int scull_open(struct inode *inode, struct file *filp)
{
    struct scull_dev *dev;
    dev = container_of(inode->i_cdev, struct scull_dev, cdev);
    filp->private_data = dev; /* diger metodlar icin */

    if ( (filp->f_flags & O_ACCMODE) == O_WRONLY) {
        scull_trim(dev); /* hatalari yoksay */
    }
    return 0;
}
\end{lstlisting}

\texttt{O\_ACCMODE} maskesi: \texttt{O\_WRONLY}/\texttt{O\_RDONLY}/
\texttt{O\_RDWR} bayrakları temiz bir bit maskesi oluşturmadığı için,
sadece erişim-modu bitlerini izole etmek amacıyla kullanılır. scull'un
\texttt{open}'ının tek gerçek işi budur: dosya \emph{sadece yazmak
için} açılmışsa cihazı sıfıra kısaltmak (normal dosyaların
\texttt{O\_APPEND} olmadan yazma-açılışında truncate edilmesiyle aynı
mantık).

\subsection{scull'un Bellek Y\"onetimi: Quantum ve Qset}

Bu bölüm zorunlu bir sürücü tekniği değil, \textbf{bir tasarım kararının}
gösterimidir: scull cihaz boyutunu bilerek sınırlamaz (felsefi olarak
keyfi sınır kötüdür; pratikte de düşük bellekli test için tüm RAM'i
tüketebilmek faydalıdır -- \texttt{cp /dev/zero /dev/scull0}).

\textbf{Veri yapısı}: her scull cihazı bir \textbf{bağlı liste}dir,
listenin her düğümü bir \texttt{scull\_qset} yapısına işaret eder:

\begin{lstlisting}
struct scull_qset {
    void **data;             /* pointer dizisi */
    struct scull_qset *next; /* bir sonraki dugum */
};
\end{lstlisting}

Terminoloji: bir işaretçinin gösterdiği bellek alanına \textbf{quantum}
denir; işaretçi dizisinin kendisine (ya da uzunluğuna)
\textbf{quantum set (qset)} denir. Öntanımlı olarak dizi 1000
işaretçi, her quantum 4000 byte'tır -- yani bir düğüm 4 milyon byte'a
kadar veri adresleyebilir.

\begin{notkutu}{Bellek fazlalığı örneği (kitaptan aynen)}
scull'a tek bir byte yazmak, o baytı tutmak için 4000 byte'lık bir
quantum \textbf{ve} 4000 (32-bit) ya da 8000 (64-bit) byte'lık bir
qset dizisi ayırır -- yani toplam 8000--12000 byte harcanır! Büyük
yazmalarda bu fazlalık önemsizdir (her 4 MB veri için sadece bir
liste elemanı). Quantum/qset boyutu \textbf{politika} meselesidir,
sürücüde sabit kodlanmamalı -- üç seviyede ayarlanabilir: derleme
zamanı (\texttt{scull.h}'deki makrolar), yükleme zamanı (modül
parametreleri), çalışma zamanı (ioctl, Bölüm 6).
\end{notkutu}

\texttt{scull\_trim()} tüm cihazın belleğini serbest bırakır (hem
\texttt{scull\_open}'ın \texttt{O\_WRONLY} truncate'inde, hem modül
kaldırılırken çağrılır):

\begin{lstlisting}
int scull_trim(struct scull_dev *dev)
{
    struct scull_qset *next, *dptr;
    int qset = dev->qset;
    int i;
    for (dptr = dev->data; dptr; dptr = next) {
        if (dptr->data) {
            for (i = 0; i < qset; i++)
                kfree(dptr->data[i]);
            kfree(dptr->data);
            dptr->data = NULL;
        }
        next = dptr->next;
        kfree(dptr);
    }
    dev->size = 0;
    dev->quantum = scull_quantum;
    dev->qset = scull_qset;
    dev->data = NULL;
    return 0;
}
\end{lstlisting}

Bu fonksiyon listeyi baştan sona gezer; her düğümde önce her bir
quantum'u, sonra dizinin kendisini, sonra düğümün kendisini serbest
bırakır (sırayla -- \texttt{next}'i düğümü silmeden önce kaydederek).

\subsection{read ve write: Tam Algoritmik Detay}

İmzalar:
\begin{lstlisting}
ssize_t read(struct file *filp, char __user *buff,
             size_t count, loff_t *offp);
ssize_t write(struct file *filp, const char __user *buff,
              size_t count, loff_t *offp);
\end{lstlisting}

\subsubsection{Neden buff'a doğrudan dokunulamaz? (3 gerekçe, kitaptan)}
\begin{enumerate}[leftmargin=1.2cm]
  \item Mimariye ve kernel ayarlarına bağlı olarak, user-space
  işaretçisi kernel modunda çalışırken \textbf{hiç geçerli
  olmayabilir}.
  \item Kullanıcı belleği \textbf{swap edilebilir} -- işaretçiye
  doğrudan dokunmak beklenmedik bir page fault tetikleyebilir; kernel
  kodu bunu doğrudan yaşayamaz (bir "oops" ile process ölür).
  \item \textbf{Güvenlik}: işaretçi kullanıcı tarafından verilir,
  hatalı ya da kötü niyetli olabilir. Körlemesine dereference etmek
  "bir user-space programının sistemin herhangi bir yerindeki belleğe
  erişmesine/üzerine yazmasına açık bir kapı" olur.
\end{enumerate}

Bu yüzden \texttt{<asm/uaccess.h>}'deki özel, mimariye-güvenli
fonksiyonlar zorunludur -- senin zaten kullandığın
\texttt{copy\_to\_user}/\texttt{copy\_from\_user}.

\begin{notkutu}{Kritik gotcha: dönüş değeri "kalan" byte sayısıdır}
\texttt{copy\_to\_user}/\texttt{copy\_from\_user}'ın dönüş değeri,
\textbf{kopyalanan} byte sayısı değil, \textbf{kopyalanamayan (kalan)}
byte sayısıdır! İşaretçi tamamen geçersizse hiç kopyalama yapılmaz;
kısmen geçersizse (örn. haritalanmamış bir sayfaya denk gelirse)
sadece bir kısmı kopyalanır -- her iki durumda da dönüş değeri "kalan"
miktardır. \textbf{scull'un kuralı}: dönüş değeri sıfır değilse
\texttt{-EFAULT} dön. Senin \texttt{ahmet.c}'in de bu kuralı doğru
uyguluyor: \texttt{if (copy\_from\_user(...)) return -EFAULT;} --
burada \texttt{if} zaten "sıfırdan farklıysa" anlamına geliyor,
doğru kullanım bu.

Ayrıca: \texttt{\_\_copy\_to\_user}/\texttt{\_\_copy\_from\_user}
(çift alt çizgili versiyonlar) adres geçerlilik kontrolünü
\textbf{atlar} -- sadece işaretçiyi önceden \texttt{access\_ok} ile
doğruladıysan kullanılmalı; doğrulama olmadan kullanmak kernel çökmesi
veya güvenlik açığına yol açabilir.
\end{notkutu}

\subsubsection{ssize\_t dönüş değeri sözleşmesi (hem read hem write için)}
\begin{itemize}[leftmargin=1.2cm]
  \item \textbf{Negatif} = hata (\texttt{<linux/errno.h>}'deki
  kodlardan biri).
  \item \textbf{\texttt{count}'a eşit} = istenen tüm veri aktarıldı
  (ideal durum).
  \item \textbf{Pozitif ama \texttt{count}'tan az} = kısmi aktarım;
  çağıran (örn. \texttt{fread}) kalanı almak için tekrar çağırmalı.
  \item \textbf{Sıfır}: \texttt{read} için = EOF (dosya sonu).
  \texttt{write} için = hiçbir şey yazılmadı ama bu \textbf{hata
  değildir} (bloklayan yazmanın tam anlamı Bölüm 6'da).
\end{itemize}

\begin{notkutu}{Gecikmeli hata bildirimi kuralı}
Eğer bir kısım veri başarıyla aktarılıp \textbf{sonra} hata oluşursa,
kural şudur: dönüş değeri başarıyla aktarılan byte sayısı olmalı, ve
\textbf{hata bir sonraki çağrıya kadar bildirilmemelidir} -- yani
sürücü bekleyen hata durumunu içeride hatırlamak zorundadır.
\end{notkutu}

scull'un tam \texttt{read} kodu (semafor çağrıları Bölüm 5'e kadar
"görmezden gel" diyerek tanıtılıyor):

\begin{lstlisting}
ssize_t scull_read(struct file *filp, char __user *buf, size_t count,
                loff_t *f_pos)
{
    struct scull_dev *dev = filp->private_data;
    struct scull_qset *dptr;
    int quantum = dev->quantum, qset = dev->qset;
    int itemsize = quantum * qset;
    int item, s_pos, q_pos, rest;
    ssize_t retval = 0;

    if (down_interruptible(&dev->sem))
        return -ERESTARTSYS;
    if (*f_pos >= dev->size)
        goto out;
    if (*f_pos + count > dev->size)
        count = dev->size - *f_pos;

    item = (long)*f_pos / itemsize;
    rest = (long)*f_pos % itemsize;
    s_pos = rest / quantum; q_pos = rest % quantum;

    dptr = scull_follow(dev, item);
    if (dptr == NULL || !dptr->data || ! dptr->data[s_pos])
        goto out; /* bosluklari doldurma */

    if (count > quantum - q_pos)
        count = quantum - q_pos;

    if (copy_to_user(buf, dptr->data[s_pos] + q_pos, count)) {
        retval = -EFAULT;
        goto out;
    }
    *f_pos += count;
    retval = count;

  out:
    up(&dev->sem);
    return retval;
}
\end{lstlisting}

Algoritma adım adım: önce \texttt{count}, \texttt{*f\_pos + count}
cihaz boyutunu aşmasın diye kırpılır (EOF'un ötesine okuma yok).
Sonra \texttt{item}/\texttt{s\_pos}/\texttt{q\_pos} -- tam sayı
bölme/mod ile -- \texttt{*f\_pos}'un hangi liste düğümüne, hangi
qset dizisi konumuna, hangi quantum-içi byte ofsetine denk geldiğini
hesaplar. \texttt{scull\_follow} listeyi gezip doğru düğümü bulur
(yoksa \texttt{NULL} döner -- "boşlukları doldurma" kuralı: eksik veri
hata değil, EOF gibi davranılır). \texttt{count} bu kez \emph{tek bir
quantum}ı aşmayacak şekilde tekrar kırpılır -- \textbf{scull'un
\texttt{read}'i bir çağrıda sadece tek bir quantum işler}, daha
fazlası isteniyorsa çağıran tekrar çağırır (stdio bunu fark etmez
bile).

\texttt{write} mantığı büyük ölçüde simetriktir, ama önemli bir
farkla: \texttt{read} "boşluk" bulursa EOF gibi davranıp durur, ama
\texttt{write} eksik olan \texttt{data} dizisini ve quantum'u
\textbf{talep üzerine (on-demand) \texttt{kmalloc} ile ayırır} --
cihaz büyüyebilir. Son olarak \texttt{dev->size}, yeni
\texttt{*f\_pos} eskisini aşıyorsa güncellenir.

\begin{ahmetkutu}
Senin \texttt{ahmet\_write}'ın \texttt{*offset}'i hiç kullanmıyor --
her yazma tüm tamponu (\texttt{buffer\_size} dahil) baştan
değiştiriyor. scull'un \texttt{write}'ı ise \texttt{*f\_pos}'u tam
olarak \texttt{read} ile aynı mantıkla kullanıyor: dosyanın
\textbf{neresine} yazıldığını hesaba katıyor. Bu, README'de zaten not
edilen eksikliğin kitaptaki "doğru" referans uygulamasıdır --
\texttt{scull\_write}'ı adım adım takip ederek kendi
\texttt{ahmet\_write}'ını \texttt{*offset} kullanacak şekilde
güncelleyebilirsin.
\end{ahmetkutu}

\subsection{scull ile Oynamak: strace \"Onerisi}

Kitap, \texttt{cp}, \texttt{dd}, yönlendirme ile test etmeyi, ve
özellikle \texttt{strace} kullanmayı önerir -- örneğin
\texttt{strace ls /dev > /dev/scull0} çalıştırırsan, \texttt{write}
çağrısının 4096 byte istediğini ama scull'un quantum sınırı yüzünden
4088 byte kabul edip kalanı için tekrar çağrıldığını görürsün. Bu
teknik Bölüm 4'te derinlemesine işlenir.

\subsection*{B\"ol\"um 3 -- H\i zl\i\ Referans}
\begin{itemize}[leftmargin=1.2cm]
  \item \texttt{dev\_t}, \texttt{MAJOR()}, \texttt{MINOR()},
  \texttt{MKDEV()} -- \texttt{<linux/types.h>}.
  \item \texttt{register\_chrdev\_region()}, \texttt{alloc\_chrdev\_region()},
  \texttt{unregister\_chrdev\_region()} -- \texttt{<linux/fs.h>}.
  \item \texttt{struct file\_operations}, \texttt{struct file},
  \texttt{struct inode} -- üç temel yapı.
  \item \texttt{cdev\_alloc()}, \texttt{cdev\_init()}, \texttt{cdev\_add()},
  \texttt{cdev\_del()} -- \texttt{<linux/cdev.h>}.
  \item \texttt{container\_of(ptr, tip, alan)} -- \texttt{<linux/kernel.h>}.
  \item \texttt{copy\_from\_user()}, \texttt{copy\_to\_user()} --
  \texttt{<asm/uaccess.h>}.
\end{itemize}

\newpage

%=================================================================
\section{Hata Ay\i klama Teknikleri}
\label{sec:bolum4}
%=================================================================

Kernel'de hata ayıklamak user space'ten yapısal olarak farklıdır: kod
tek bir process'e bağlı değildir, hatalı kod tüm sistemi çökertebilir,
ve çökme "kanıtı" genelde kaybolur. Bu bölüm, sırasıyla \emph{yazdırarak},
\emph{sorgulayarak}, \emph{izleyerek} ve son çare olarak
\emph{interaktif debugger ile} hata ayıklamayı ele alır.

\subsection{printk: Log Seviyeleri}

\texttt{printk}'in \texttt{printf}'ten temel farkı: format string'in
başına bir \textbf{öncelik (loglevel)} dizgesi eklenebilir. Sekiz
seviye vardır (\texttt{<linux/kernel.h>}), \textbf{azalan önem
sırasıyla}:

\begin{center}
\begin{tabular}{ll}
\texttt{KERN\_EMERG} & sistem kullanılamaz duruma gelmeden hemen önce \\
\texttt{KERN\_ALERT} & derhal müdahale gerektiren durum \\
\texttt{KERN\_CRIT} & kritik durumlar, ciddi donanım/yazılım hatası \\
\texttt{KERN\_ERR} & hata durumu (sürücüler donanım sorunları için kullanır) \\
\texttt{KERN\_WARNING} & ciddi olmayan ama sorunlu durumlar \\
\texttt{KERN\_NOTICE} & normal ama dikkate değer durumlar \\
\texttt{KERN\_INFO} & bilgilendirme (örn. açılışta bulunan donanım) \\
\texttt{KERN\_DEBUG} & hata ayıklama mesajları \\
\end{tabular}
\end{center}

Loglevel belirtilmezse öntanımlı \texttt{KERN\_WARNING} kullanılır.
Mesaj, kendi öncelik sayısı \texttt{console\_loglevel} değişkeninden
\textbf{küçükse} konsola da basılır; \texttt{klogd}/\texttt{syslogd}
çalışıyorsa mesaj bağımsız olarak \texttt{/var/log/messages}'a da
yazılır. \texttt{console\_loglevel} değiştirme yöntemleri:
\texttt{klogd -c <seviye>}, ya da doğrudan
\texttt{/proc/sys/kernel/printk} dosyasına yazmak (\texttt{echo 8 >
/proc/sys/kernel/printk} $\to$ debug dahil her şeyi konsola bas).

\begin{ahmetkutu}
Senin \texttt{ahmet.c}'in her yerde \texttt{pr\_info} kullanıyor --
bu, \texttt{printk(KERN\_INFO ...)} için kısayol bir makrodur (aynı
aile: \texttt{pr\_err}, \texttt{pr\_warn}, \texttt{pr\_debug}...).
\texttt{dmesg} ile gördüğün çıktı tam olarak bu mekanizmanın
sonucudur.
\end{ahmetkutu}

\subsubsection{printk nas\i l saklan\i r: dairesel tampon (circular buffer)}

\texttt{printk}, verisini sabit boyutlu bir \textbf{dairesel tampona}
yazar (\texttt{\_\_LOG\_BUF\_LEN}, 4KB--1MB arası, derleme zamanı
seçimi). Tampon dolarsa \textbf{en eski veri üzerine yazılır} --
bilinçli bir tercih: bu sayede bir loglama süreci (process) olmadan
da sistem çalışabilir, bellek de israf edilmez. \texttt{dmesg} bu
tamponu \textbf{tüketmeden} (okuyup silmeden) dump eder; klogd'un
kullandığı \texttt{/proc/kmsg} okuması ise tamponu tüketir (FIFO
gibi davranır -- ikinci bir okuyucu varsa çakışma olur).

\subsubsection{printk\_ratelimit: log ta\c{s}mas\i n\i\ \"onlemek}

Dikkatsiz \texttt{printk} kullanımı binlerce mesaj üretip konsolu
boğabilir; \textbf{yavaş konsol cihazlarında} (örn. seri port) sistemi
\textbf{yavaşlatabilir hatta kilitleyebilir} (kesme kullanılamayan bir
seri konsolu sürmek için). \textbf{Kural: üretim kodu normal
çalışmada asla yazdırmamalı, sadece istisnai durumlarda.} Tekrarlayan
hata durumları için kernel bir hız sınırlayıcı sağlar:

\begin{lstlisting}
if (printk_ratelimit())
    printk(KERN_NOTICE "The printer is still on fire\n");
\end{lstlisting}

Ayarlanabilir eşikler: \texttt{/proc/sys/kernel/printk\_ratelimit}
(saniye), \texttt{/proc/sys/kernel/printk\_ratelimit\_burst} (eşiğe
kadar kabul edilen mesaj sayısı).

\subsection{Ko\c{s}ullu Debug Mesajlar\i : PDEBUG Kal\i b\i}

Kitabın scull kaynağında kullandığı, derleme-zamanı açılıp
kapanabilen debug makrosu:

\begin{lstlisting}
#undef PDEBUG
#ifdef SCULL_DEBUG
#  ifdef __KERNEL__
#    define PDEBUG(fmt, args...) printk( KERN_DEBUG "scull: " fmt, ## args)
#  else
#    define PDEBUG(fmt, args...) fprintf(stderr, fmt, ## args)
#  endif
#else
#  define PDEBUG(fmt, args...) /* debug yok: hicbir sey */
#endif
\end{lstlisting}

Bu, gcc'nin değişken-argümanlı makro genişletmesini kullanır. Kernel
alanında \texttt{printk}, kullanıcı alanında (aynı kodu test programı
olarak derlersen) \texttt{fprintf(stderr,...)} olur -- \textbf{aynı
debug makrosu iki dünyada da çalışır}. Makefile'da
\texttt{-DSCULL\_DEBUG} bayrağıyla açılır/kapanır -- yani debug
mesajlarını değiştirmek \textbf{yeniden derleme} gerektirir (buna
karşın çalışma-zamanı \texttt{if} koşuluyla yapılsaydı her zaman ufak
bir performans maliyeti olurdu, ama yeniden derleme gerekmezdi --
kitap bu ödünleşimi (trade-off) açıkça belirtir).

\subsection{/proc Dosya Sistemi ile Sorgulama}

Sürekli \texttt{printk} basmak yerine, durumu \textbf{talep üzerine}
(\texttt{ps}, \texttt{netstat} gibi) sorgulamak daha az müdahaleci bir
yöntemdir. \texttt{/proc}, her dosyası bir kernel fonksiyonuna bağlı,
yazılım tarafından anlık üretilen sanal bir dosya sistemidir.

\begin{notkutu}{Kitabın uyarısı: /proc kullanımı art\i k \"onerilmiyor}
Kitap açıkça belirtiyor: kernel geliştiricileri \texttt{/proc}'a yeni
dosya eklemeyi \textbf{caydırıyor} -- orijinal amacının (process bilgisi)
dışına taşıp "kontrolsüz bir çöplük" haline geldiği düşünülüyor. Modern
önerilen yöntem \textbf{sysfs}'tir (Bölüm 14). Kitap yine de
\texttt{/proc}'u anlatır çünkü hata ayıklama amaçlı basit kullanım
için yeterlidir ve daha kolaydır.
\end{notkutu}

Basit salt-okunur bir \texttt{/proc} girdisi için (2.6.10-dönemi API):

\begin{lstlisting}
int (*read_proc)(char *page, char **start, off_t offset, int count,
                  int *eof, void *data);

struct proc_dir_entry *create_proc_read_entry(const char *name,
        mode_t mode, struct proc_dir_entry *base,
        read_proc_t *read_proc, void *data);
\end{lstlisting}

Kernel, sürücünün veri yazması için \texttt{PAGE\_SIZE} boyutunda bir
tampon (\texttt{page}) ayırır; \texttt{count}/\texttt{offset} normal
\texttt{read} ile aynı anlama gelir; \texttt{eof} sürücünün "veri
bitti" diye işaretlediği bir \texttt{int} işaretçisidir. Modül
kaldırılırken \texttt{remove\_proc\_entry()} ile girdi
\textbf{mutlaka} silinmelidir.

\begin{notkutu}{İki gotcha: proc girdileri modülün refcount'una bağlı değil}
(1) \texttt{/proc} girdilerinin bir "sahibi" (owner) yoktur -- yani
bir process dosyayı açıkken (\texttt{sleep 100 < /proc/dosyam})
modül kaldırılırsa \textbf{çökme riski} vardır. (2) Kernel aynı
isimde iki \texttt{/proc} girdisi oluşturulmasını \textbf{kontrol
etmez} -- kitabın kendi ifadesiyle bu, "sınıflarda olduğu bilinen"
bir hata kaynağıdır.
\end{notkutu}

Büyük çıktılar için modern/tercih edilen yaklaşım \textbf{seq\_file}
arayüzüdür (\texttt{<linux/seq\_file.h>}) -- \texttt{start}/
\texttt{next}/\texttt{stop}/\texttt{show} adlı dört yineleyici (iterator)
fonksiyonu uygulanır, çıktı \texttt{seq\_printf} ile üretilir (ham
\texttt{printk} ile değil). Kitap: "birden fazla küçük miktarda çıktı
üreten her dosya için seq\_file kullanılmalı."

\subsection{ioctl ile Hata Ay\i klama}

\texttt{/proc}'a alternatif: sürücüye özel bir debug \texttt{ioctl}
komutu ekleyip iç veri yapılarını user space'e kopyalamak. Avantajı:
\texttt{/proc} dosyaları dizin listelemesinde \textbf{herkese görünür}
(meraklı gözlerin dikkatini çeker), oysa belgelenmemiş bir ioctl
komutu fark edilmesi zor ama gerektiğinde hâlâ oradadır. Dezavantajı:
kullanmak için ayrı bir user-space programı yazman gerekir.

\subsection{strace: Sistem \c{C}a\u{g}r\i lar\i n\i\ D\i\c{s}ar\i dan \.{I}zlemek}

\texttt{strace}, bir user-space programının yaptığı \textbf{tüm sistem
çağrılarını} -- isim, argümanlar, dönüş değeri -- sembolik biçimde
gösterir. Kernel'in kendisinden bilgi aldığı için, izlenen programın
\texttt{-g} ile derlenmiş olması gerekmez.

\begin{notkutu}{Kitabın gerçek örneği: scull'un quantum sınırı strace'te nasıl görünür}
\texttt{strace ls /dev > /dev/scull0} çalıştırıldığında:
\texttt{write(1, ..., 4096) = 4088} sonra kalan 8 byte için ikinci bir
\texttt{write} çağrısı görülür -- çünkü scull'un \texttt{write}'ı bir
seferde sadece tek bir quantum (4000 byte civarı) kabul eder. Bu,
\texttt{strace}'in bir sürücünün iç davranışını nasıl \textbf{dışarıdan
gözlemlenebilir} kıldığının somut bir örneğidir.
\end{notkutu}

\subsection{Oops Mesajlar\i n\i\ Okumak}

Çoğu kernel hatası \texttt{NULL} işaretçi dereference'ı ya da geçersiz
bir işaretçi değeri şeklinde ortaya çıkar $\to$ bir \textbf{oops}
mesajı. \textbf{Oops $\neq$ panic}: Linux dayanıklıdır, bir oops
genelde sadece o anki process'i öldürür, sistem çalışmaya devam eder
(panic ise tüm sistemi durdurur, çok daha nadirdir).

Örnek (kitaptan, \texttt{faulty} modülünün \texttt{*(int*)0 = 0;}
satırından):

\begin{lstlisting}
Unable to handle kernel NULL pointer dereference at virtual address 00000000
Oops: 0002 [#1]
EIP:    0060:[<d083a064>]    Not tainted
EIP is at faulty_write+0x4/0x10 [faulty]
...
Call Trace:
 [<c0150558>] vfs_write+0xb8/0x130
 [<c0150682>] sys_write+0x42/0x70
\end{lstlisting}

Okuma stratejisi (kitabın kendi sırası):
\begin{enumerate}[leftmargin=1.2cm]
  \item \textbf{Önce \texttt{EIP is at ...} satırına bak} -- hatanın
  tam olarak nerede olduğunu (fonksiyon adı + modül adı + fonksiyon
  içindeki offset) sembolik olarak gösterir. Bu örnekte:
  \texttt{faulty\_write} fonksiyonu, \texttt{faulty} modülü, fonksiyonun
  başından 4 (hex) byte içeride, fonksiyon toplam 0x10 byte uzunlukta.
  \item Yeterli değilse, \textbf{Call Trace} (çağrı yığını) bloğuna
  bak -- oraya nasıl gelindiğini gösterir.
  \item \textbf{"Tainted" bayrağı}na dikkat et -- GPL-olmayan/güvenilmez
  bir modül yüklenmişse "Tainted" görünür, bu da hata raporunun ciddiye
  alınma olasılığını etkiler.
\end{enumerate}

\begin{notkutu}{Sembolik traceback için gereken kernel ayarı}
Call Trace'in \textbf{sembolik} (fonksiyon adlarıyla) görünmesi için
kernel \texttt{CONFIG\_KALLSYMS} ile derlenmiş olmalı -- yoksa sadece
ham hex adresler görürsün, hata ayıklama çok daha zorlaşır.
\end{notkutu}

\begin{notkutu}{Slab poison değerleri: \texttt{0xa5} ve \texttt{0x6b}}
\texttt{CONFIG\_DEBUG\_SLAB} açıkken, kernel'in bellek ayırıcısı
tahsis edilen her byte'ı çağırana vermeden önce \texttt{0xa5} deseniyle,
serbest bırakılan belleği ise \texttt{0x6b} deseniyle doldurur. Bir
oops mesajında ya da sürücü çıktısında tekrar eden \texttt{0xa5} veya
\texttt{0x6b} (ya da \texttt{0xa5a5a5a5} gibi bir adres) görürsen, bu
\textbf{başlatılmamış ya da serbest bırakılmış belleğin kullanıldığının}
klasik bir işaretidir.
\end{notkutu}

\subsection{Sistem Kilitlenmeleri: Magic SysRq}

Bazı hatalar oops bile vermeden sistemi \textbf{tamamen kilitler}.
Kod sonsuz bir döngüye girerse, kernel zamanlamayı (scheduling)
durdurabilir -- sistem Ctrl-Alt-Del dahil hiçbir şeye yanıt vermez.
Önleyici teknik: uzun döngülerin içine stratejik noktalarda
\texttt{schedule()} çağrıları eklemek (\textbf{ama asla bir spinlock
tutarken \texttt{schedule()} çağırma}).

\textbf{Magic SysRq} tuşu (Alt+SysRq+komut), çoğu kilitlenmede
vazgeçilmez bir araçtır -- \texttt{CONFIG\_MAGIC\_SYSRQ} ile
etkinleştirilmelidir (çoğu dağıtımda güvenlik nedeniyle
\textbf{varsayılan olarak kapalıdır}). Faydalı komutlar:

\begin{itemize}[leftmargin=1.2cm]
  \item \texttt{r} -- klavyeyi raw moddan çıkarır (çökmüş bir X
  sunucusunun bıraktığı garip klavye durumunu düzeltir).
  \item \texttt{s} -- tüm diskleri acil senkronize eder (sync).
  \item \texttt{u} -- tüm diskleri salt-okunur yeniden bağlar (umount);
  genelde \texttt{s}'den hemen sonra, dosya sistemi kontrolü süresini
  kısaltmak için.
  \item \texttt{b} -- anında yeniden başlatır (önce \texttt{s} ve
  \texttt{u} çalıştırılmalı!).
  \item \texttt{p} -- işlemci register bilgisini yazdırır.
  \item \texttt{t} -- güncel görev (task) listesini yazdırır.
\end{itemize}

\subsection{Debugger'lar: gdb, kdb, kgdb, UML}

\begin{itemize}[leftmargin=1.2cm]
  \item \textbf{gdb + /proc/kcore}: \texttt{gdb /usr/src/linux/vmlinux
  /proc/kcore} -- \texttt{/proc/kcore}, tüm kernel adres alanını temsil
  eden, okundukça \emph{üretilen} özel bir dosyadır. \textbf{Sınırlama}:
  gdb veri okur ama \textbf{önbelleğe alır} -- \texttt{p jiffies} gibi
  değişen bir değeri tekrar sorgularsan \textbf{eski (bayat) veri}
  görebilirsin; \texttt{core-file /proc/kcore} komutunu tekrar
  vererek önbelleği temizlemen gerekir. Ayrıca gdb bu modda
  \textbf{breakpoint koyamaz, veri değiştiremez, tek adım
  ilerleyemez} -- sadece okuma amaçlıdır. \texttt{CONFIG\_DEBUG\_INFO}
  gerektirir.
  \item \textbf{kdb}: resmi kernel'e dahil değildir (Linus Torvalds
  interaktif debugger'lara felsefi olarak karşıdır -- "semptomu
  yamalamayı, kök nedeni bulmayı değil teşvik ediyor" endişesiyle),
  ayrı bir yama olarak dağıtılır, sadece x86'da çalışır.
  \item \textbf{kgdb}: gerçek breakpoint/tek-adım desteği için iki
  ayrı makine kullanır -- test edilen kernel bir "kurban" makinede,
  gdb'nin kendisi \textbf{stabil bir masaüstünde} çalışır, aralarında
  genelde bir \textbf{seri kablo} bağlantısı vardır.
  \item \textbf{User-Mode Linux (UML)}: kernelin kendisini normal bir
  user-space process'i olarak çalıştıran ayrı bir port. Avantajı:
  hatalı kernel gerçek sistemi bozamaz. \textbf{Kritik kısıt}: UML
  kernelinin gerçek donanıma \textbf{erişimi yoktur} -- donanımla
  konuşan sürücü kodunu debug etmek için kullanılamaz.
\end{itemize}

\subsection*{B\"ol\"um 4 -- \"On Plana \c{C}\i kan Uyar\i lar}
\begin{itemize}[leftmargin=1.2cm]
  \item Debug config seçenekleri (\texttt{CONFIG\_DEBUG\_*}) ekstra
  çıktı üretir ve performansı düşürür -- bu yüzden üretim/dağıtım
  kernellerinde bilinçli olarak kapalıdır.
  \item \texttt{CONFIG\_INPUT\_EVBUG} \textbf{her tuş vuruşunu, şifreler
  dahil} loglar -- açık bir güvenlik uyarısı.
  \item \texttt{/proc} girdisi kaldırmayı unutmak, modül kaldırılırken
  beklenmedik çökmelere yol açar.
  \item \texttt{schedule()} çağrıları reentrant çağrılara kapı açar --
  asla bir spinlock tutarken çağrılmamalı.
\end{itemize}

\newpage

%=================================================================
\section{E\c{s}zamanl\i l\i k ve Race Condition'lar}
\label{sec:bolum5}
%=================================================================

\begin{notkutu}{Bu bölüm neden senin için en kritik olabilir}
Senin \texttt{ahmet.c}'inde \textbf{şu an hiçbir kilitleme yok} --
global \texttt{buffer} ve \texttt{buffer\_size} değişkenleri
korumasız. Bu bölüm tam olarak bunun neden tehlikeli olduğunu ve nasıl
düzeltileceğini anlatıyor.
\end{notkutu}

\subsection{scull'daki Somut Race Condition}

Kitap, scull'un \texttt{write} kodundan gerçek bir parça üzerinden
race condition'ı gösterir:

\begin{lstlisting}
if (!dptr->data[s_pos]) {
    dptr->data[s_pos] = kmalloc(quantum, GFP_KERNEL);
    if (!dptr->data[s_pos])
        goto out;
}
\end{lstlisting}

\textbf{Senaryo}: İki process, A ve B, \textbf{aynı scull cihazının aynı
ofsetine} aynı anda yazıyor.
\begin{enumerate}[leftmargin=1.2cm]
  \item Her ikisi de \texttt{if (!dptr->data[s\_pos])} testine
  \textbf{aynı anda} ulaşır.
  \item İşaretçi \texttt{NULL} olduğu için \textbf{ikisi de} bellek
  ayırmaya karar verir ve sonucu aynı yere atar.
  \item İki atama da aynı konuma yazdığı için, sadece \textbf{sonuncusu}
  "kazanır."
  \item scull, \textbf{ilk (kaybeden) process'in ayırdığı belleği
  tamamen unutur} -- bu bellek asla serbest bırakılmaz $\to$
  \textbf{bellek sızıntısı}.
\end{enumerate}

Kitabın tanımı: \textbf{"Race condition, paylaşılan veriye kontrolsüz
erişimin bir sonucudur. Yanlış erişim örüntüsü gerçekleştiğinde,
beklenmedik bir şey olur."} Ve önemli bir uyarı: \textbf{"Programcılar
race condition'ları son derece düşük olasılıklı olaylar olarak
görmezden gelmeye meyillidir. Ama bilgi işlem dünyasında, milyonda bir
olay her birkaç saniyede bir gerçekleşebilir, ve sonuçları ciddi
olabilir."}

\subsection{Genel Kural: Payla\c{s}\i lan Kaynak = Y\"onetim Zorunlulu\u{g}u}

Kitabın "sert kural"ı (tam ifadeyle): \textbf{"bir donanım ya da
yazılım kaynağı tek bir çalışma iş parçacığının (thread of execution)
ötesinde paylaşıldığında, ve bir iş parçacığının o kaynağın tutarsız
bir görünümüyle karşılaşma olasılığı varsa, o kaynağa erişimi açıkça
yönetmek zorundasın."}

Eşzamanlılığın kaynakları (kitabın listesi):
\begin{enumerate}[leftmargin=1.2cm]
  \item Birden fazla user-space process'in sürücü koduna "şaşırtıcı
  kombinasyonlarda" erişmesi.
  \item SMP sistemlerde farklı işlemcilerin sürücü kodunu \textbf{gerçekten
  aynı anda} çalıştırması.
  \item \textbf{Kernel preemption} -- 2.6 kernel kodu önceliklenebilir
  (preemptible); sürücü kodun herhangi bir anda işlemciyi kaybedebilir,
  ve yerine geçen process de aynı sürücü kodunda olabilir.
  \item Kesmeler (interrupt) -- asenkron olarak sürücü kodunun eşzamanlı
  çalışmasına neden olur.
  \item Ertelenmiş çalıştırma mekanizmaları (workqueue, tasklet, timer).
  \item Takılıp-çıkarılabilir (hot-pluggable) cihazlar -- cihaz tam
  onunla uğraşırken ortadan kaybolabilir.
\end{enumerate}

\subsection{Semaphore ve Mutex}

\textbf{Kritik bölge (critical section)}: "aynı anda sadece bir iş
parçacığı tarafından çalıştırılabilecek kod."

\begin{lstlisting}
DECLARE_MUTEX(name);          /* 1'e ilklendirilmis */
DECLARE_MUTEX_LOCKED(name);   /* 0'a ilklendirilmis (kilitli baslar) */
void init_MUTEX(struct semaphore *sem);
void init_MUTEX_LOCKED(struct semaphore *sem);

void down(struct semaphore *sem);              /* kesilemez uyku */
int  down_interruptible(struct semaphore *sem); /* sinyal ile kesilebilir */
int  down_trylock(struct semaphore *sem);       /* hic uyumaz */
void up(struct semaphore *sem);
\end{lstlisting}

\begin{notkutu}{down\_interruptible neredeyse her zaman doğru seçimdir}
\texttt{down()} (kesilemez) kullanmak "öldürülemeyen process'ler
yaratmanın iyi bir yolu" (ps çıktısındaki korkulan \texttt{D} durumu) --
kullanıcıları rahatsız eder. \textbf{Kural}: \texttt{down\_interruptible}
sıfırdan farklı dönerse, semafor \textbf{tutulmuyor demektir} --
dönüş değeri her zaman kontrol edilmeli, ve tipik olarak
\texttt{-ERESTARTSYS} döndürülmelidir.
\end{notkutu}

\begin{notkutu}{ERESTARTSYS kuralı}
\texttt{-ERESTARTSYS} döndürüyorsan, önce kullanıcıya görünür herhangi
bir değişikliği geri almalısın (ki syscall yeniden başladığında doğru
çalışsın); geri alamıyorsan bunun yerine \texttt{-EINTR} döndürmelisin.
\end{notkutu}

\subsubsection{scull'da semafor kullan\i m\i}

\begin{lstlisting}
struct scull_dev {
    ...
    struct semaphore sem;  /* karsilikli dislama semaforu */
    struct cdev cdev;
};

for (i = 0; i < scull_nr_devs; i++) {
    scull_devices[i].quantum = scull_quantum;
    scull_devices[i].qset = scull_qset;
    init_MUTEX(&scull_devices[i].sem);
    scull_setup_cdev(&scull_devices[i], i);
}
\end{lstlisting}

\begin{notkutu}{Kritik sıralama kuralı}
Semafor, \textbf{cihaz sistemin geri kalanına açılmadan önce}
ilklendirilmeli -- yani \texttt{init\_MUTEX}, \texttt{scull\_setup\_cdev}
(içinde \texttt{cdev\_add} çağıran) \textbf{çağrılmadan önce}
gelmelidir. Sıra tersine çevrilirse, "semaforun hazır olmadan önce
erişilebileceği bir race condition yaratılır." Bu aynen Bölüm 2'de
gördüğün "bir yeteneği tamamen hazır olmadan kaydetme" kuralının somut
bir örneğidir.
\end{notkutu}

\textbf{Neden global tek bir semafor değil, cihaz-başına semafor?}
Kitap: "çeşitli scull cihazları hiçbir kaynağı ortak kullanmıyor...
bir process farklı bir scull cihazıyla çalışırken diğerini beklemeye
zorlamak için hiçbir gerekçe yok" -- paralelliği artırır.

\subsection{Reader/Writer Semaphore (rwsem)}

Çoğu görev salt-okunur/yazma olarak ayrışır; \texttt{struct
rw\_semaphore} birden fazla eşzamanlı okuyucuya (yazar yokken) izin
verir:

\begin{lstlisting}
void down_read(struct rw_semaphore *sem);
int  down_read_trylock(struct rw_semaphore *sem);  /* NOT: basari=sifirdan FARKLI, tersi degil! */
void up_read(struct rw_semaphore *sem);
void down_write(struct rw_semaphore *sem);
void up_write(struct rw_semaphore *sem);
\end{lstlisting}

\begin{notkutu}{İnce ama önemli tersine dönüş}
\texttt{down\_read\_trylock}/\texttt{down\_write\_trylock} çoğu kernel
fonksiyonunun aksine \textbf{sıfırdan farklı = başarı, sıfır = başarısız}
kuralını kullanır -- kitap bunu açıkça "çoğu kernel fonksiyonunun
tersi bir kural" diye işaretliyor.
\end{notkutu}

\textbf{rwsem'lerde yazarlar önceliklidir} -- bir yazar beklemeye
başladığında yeni okuyucular kabul edilmez, bu da \textbf{okuyucu
açlığına} (reader starvation) yol açabilir. Bu yüzden rwsem'ler sadece
yazma erişiminin nadir ve kısa olduğu durumlarda tercih edilmeli.

\subsection{Completion: "\.{I}\c{s} Bitti" Bildirimi}

\textbf{Yanlış yaklaşım} (kitabın açıkça uyardığı): bir işi başlatıp
kilitli bir semaforda \texttt{down()} ile beklemek, işi yapan taraf
\texttt{up()} çağırınca devam etmek. Bu yanlış çünkü: (1) semaforlar
"müsait" (uncontended) hızlı yol için optimize edilmiştir, oysa burada
çağıran neredeyse \textbf{her zaman} bekleyecektir -- tam tersi
senaryo; (2) semafor bir \textbf{stack (otomatik) değişkeni} olarak
tanımlanmışsa, bekleyen taraf fonksiyonundan dönüp semaforu yok
ederken, sinyal veren taraf hâlâ ona dokunuyor olabilir -- ciddi bir
race condition.

Bu yüzden \textbf{completion} arayüzü eklenmiştir:

\begin{lstlisting}
DECLARE_COMPLETION(my_completion);
/* veya dinamik: */
struct completion my_completion;
init_completion(&my_completion);

void wait_for_completion(struct completion *c);  /* kesilemez bekleme! */
void complete(struct completion *c);      /* sadece TEK bekleyeni uyandirir */
void complete_all(struct completion *c);  /* TUM bekleyenleri uyandirir */
\end{lstlisting}

\begin{notkutu}{wait\_for\_completion kesilemez}
\texttt{wait\_for\_completion} \textbf{kesilemez} (uninterruptible)
bir bekleme yapar -- eğer hiç kimse \texttt{complete()} çağırmazsa,
sonuç \textbf{öldürülemeyen bir process}tir. Bu yüzden bu API'yi
kullanırken, birinin gerçekten \texttt{complete} çağıracağından emin
olmalısın.
\end{notkutu}

\subsection{Spinlock: Uyuyamayan Kod \.{I}\c{c}in Kilit}

Semaforun aksine spinlock \textbf{uyku gerektirmez/desteklemez} --
kesme işleyicileri gibi uyuyamayan bağlamlarda kullanılır. Kilidi
alamayan kod, müsait olana kadar sıkı bir döngüde bekler ("spin
eder").

\begin{notkutu}{Çekirdek kural: spinlock tutarken asla uyuma}
"Bir spinlock tutan herhangi bir kod \textbf{atomik} olmak zorundadır.
Uyuyamaz; hatta kesmelere hizmet etmek dışında herhangi bir nedenle
işlemciyi bırakamaz (bazen o bile mümkün olmaz)." Neden? Eğer kilidi
tutan kod uyursa (örn. \texttt{copy\_from\_user} ya da
\texttt{kmalloc(GFP\_KERNEL)} gibi \textbf{gizlice uyuyabilen}
fonksiyonlar çağrılırsa), kilit süresiz tutulabilir -- diğer taraflar
sonsuza kadar bekler, \textbf{sistem kilitlenebilir}.
\end{notkutu}

\begin{notkutu}{Aynı-CPU kesme kilitlenmesi (deadlock) senaryosu}
Sürücü bir spinlock'u tutuyor; tutarken cihaz bir kesme (interrupt)
üretiyor, kesme işleyicisi de \textbf{aynı kilidi} istiyor. Eğer
kesme, kilidi tutan işlemcinin \textbf{aynısında} gerçekleşirse:
kesilmeyen (non-interrupt) kod artık çalışıp kilidi serbest bırakma
şansı bulamaz -- o işlemci \textbf{sonsuza kadar döner}. Çözüm: kilidi
tutarken \textbf{o işlemcideki kesmeleri devre dışı bırakan} varyantları
kullanmak.
\end{notkutu}

\begin{lstlisting}
spinlock_t my_lock = SPIN_LOCK_UNLOCKED;   /* derleme-zamani */
spin_lock_init(spinlock_t *lock);          /* calisma-zamani */

void spin_lock(spinlock_t *lock);
void spin_lock_irqsave(spinlock_t *lock, unsigned long flags);
void spin_lock_irq(spinlock_t *lock);
void spin_lock_bh(spinlock_t *lock);       /* sadece softirq/tasklet'i kapatir */

void spin_unlock(spinlock_t *lock);
void spin_unlock_irqrestore(spinlock_t *lock, unsigned long flags);
void spin_unlock_irq(spinlock_t *lock);
void spin_unlock_bh(spinlock_t *lock);
\end{lstlisting}

\begin{notkutu}{Hangi varyant, ne zaman?}
Kilit bir \textbf{donanım kesme işleyicisi} tarafından da alınabiliyorsa
$\to$ \texttt{\_irqsave}/\texttt{\_irq} formlarından biri
\textbf{zorunlu} (yoksa sistem er geç kilitlenebilir). Kilit sadece bir
tasklet/softirq bağlamından da alınıyorsa (donanım kesmesinden değil)
$\to$ \texttt{\_bh} yeterli, donanım kesmeleri açık kalabilir.
\texttt{spin\_lock\_irqsave}/\texttt{spin\_unlock\_irqrestore}
\textbf{aynı fonksiyonda, aynı \texttt{flags} değişkeniyle}
eşleştirilmeli -- aksi halde bazı mimarilerde kod bozulabilir.
\end{notkutu}

\textbf{Kilit tutma süresi kuralı}: spinlock'lar \textbf{mümkün olan
en kısa süre} tutulmalı -- uzun tutulan bir kilit, diğer işlemcilerin
bekleme süresini ve zamanlama gecikmesini (scheduling latency) doğrudan
kötüleştirir.

\subsection{Kilitleme Tuzaklar\i}

\begin{itemize}[leftmargin=1.2cm]
  \item \textbf{Self-deadlock}: bir fonksiyon bir kilidi alıp, sonra
  \textbf{aynı kilidi} tekrar almaya çalışan başka bir fonksiyonu
  çağırırsa, kod kilitlenir -- ne semaphore ne spinlock, bir kilit
  sahibinin kilidi ikinci kez almasına izin vermez, basitçe asılı
  kalır.
  \item \textbf{Sıralama (ordering) deadlock'u}: iki kilit
  (\texttt{Lock1}, \texttt{Lock2}) varsa ve bir thread önce
  \texttt{Lock1} sonra \texttt{Lock2} alırken, başka bir thread önce
  \texttt{Lock2} sonra \texttt{Lock1} almaya çalışırsa \textbf{ikisi
  de kilitlenir}. \textbf{Çözüm}: birden fazla kilit alınacaksa,
  \textbf{her zaman aynı sırayla} alınmalı. Genel kural: yerel
  (kendi) kilidini, kernelin daha merkezi bir kilidinden \textbf{önce}
  al; semaforları spinlock'lardan \textbf{önce} al (bir spinlock
  tutarken \texttt{down} çağırmak -- yani uyuyabilecek bir şey
  çağırmak -- "ciddi bir hatadır").
  \item \textbf{Büyük Kernel Kilidi (BKL)}: Linux'un ilk SMP-uyumlu
  kerneli (2.0) tek bir spinlock'a sahipti -- tüm kerneli tek bir
  dev kritik bölgeye çeviriyordu. 2.6'da hâlâ (çok azaltılmış biçimde)
  var, ama kitap açıkça uyarıyor: \textbf{"yeni kodda onu kullanmayı
  aklından bile geçirme."}
  \item \textbf{İnce mi kaba mı taneli kilitleme}: kitabın tavsiyesi
  sürücü yazarları için nettir: \textbf{"gerçek bir çekişme (contention)
  sorunu olacağına dair somut bir sebep olmadıkça, görece kaba taneli
  kilitlemeyle başla. Erken optimizasyon dürtüsüne direnç göster;
  gerçek performans kısıtlamaları genelde beklenmedik yerlerde ortaya
  çıkar."}
\end{itemize}

\subsection{Kilitlemeye Alternatifler}

\begin{itemize}[leftmargin=1.2cm]
  \item \textbf{Dairesel tampon (lock-free)}: tek yazar/tek okuyucu
  senaryosunda, yazar veriyi \textbf{önce} tamponun içine yazıp
  \textbf{sonra} yazma-indeksini güncellerse, okuyucu her zaman
  tutarlı bir görünüm görür -- \textbf{hiç kilide gerek kalmadan}.
  2.6.10'dan itibaren kernel'de hazır bir uygulaması var:
  \texttt{<linux/kfifo.h>}.
  \item \textbf{atomic\_t}: tek bir tamsayı üzerinde atomik işlemler
  (\texttt{<asm/atomic.h>}). \textbf{24 bitten fazlasına güvenme}
  (bazı mimarilerin sınırlaması).
  \begin{lstlisting}
atomic_t v = ATOMIC_INIT(0);
atomic_set(&v, i); atomic_read(&v);
atomic_inc(&v); atomic_dec(&v);
int atomic_dec_and_test(atomic_t *v); /* sonuc 0 ise true doner */
  \end{lstlisting}
  \begin{notkutu}{atomic\_t sadece TEK değişkeni korur}
  İki ayrı \texttt{atomic\_t} üzerinde art arda işlem yaparsan
  (\texttt{atomic\_sub(...,\ \&a); atomic\_add(...,\ \&b);}), ikisi
  arasında bir pencere vardır -- "a" değişmiş ama "b" henüz
  değişmemiş bir an. Bu ara durum başka bir threadi etkileyebiliyorsa,
  \textbf{yine de kilit gerekir} -- atomic\_t bileşik (çok değişkenli)
  işlemleri atomik yapmaz.
  \end{notkutu}
  \item \textbf{Bit işlemleri}: \texttt{set\_bit}, \texttt{clear\_bit},
  \texttt{test\_and\_set\_bit} vb. (\texttt{<asm/bitops.h>}) atomik bit
  erişimi sağlar. Kitap bit-tabanlı elle kilitlemeyi \textbf{caydırır}:
  "yeni kodda spinlock kullanmak çok daha iyidir; spinlock'lar iyi
  test edilmiştir, kesme/preemption gibi konuları halleder, ve kodunu
  okuyan başkalarının ne yaptığını anlamak için uğraşmasına gerek
  kalmaz."
  \item \textbf{RCU (Read-Copy-Update)}: okumanın sık, yazmanın nadir
  olduğu durumlar için ileri düzey, kilitsiz bir teknik -- veri sadece
  \textbf{işaretçi üzerinden} erişilmeli. Yazar değişikliği önce bir
  \textbf{kopya} üzerinde yapar, sonra işaretçiyi tek bir atomik
  adımda yeni kopyaya yönlendirir; eski kopya, tüm işlemcilerin en az
  bir kez zamanlanmasından sonra (artık kimse referans tutmuyor
  garantisiyle) \texttt{call\_rcu()} ile serbest bırakılır. Gerçek
  kullanım örneği: ağ yönlendirme tabloları.
\end{itemize}

\begin{notkutu}{Kitapta AÇIKÇA yer almayan bir modern API: struct mutex}
Bu bölümün kitaptaki hali \textbf{sadece} eski semaphore-tabanlı
\texttt{DECLARE\_MUTEX}/\texttt{down}/\texttt{up} API'sini "mutex"
olarak anlatır. Daha sonraki kernel sürümlerinde eklenen hafif
\texttt{struct mutex} API'si (\texttt{mutex\_init},
\texttt{mutex\_lock}, \texttt{mutex\_unlock}) bu kitapta \textbf{hiç
geçmiyor} -- bu, kitabın 2005 tarihli olmasından kaynaklanan bir
eksiklik, sana ait bir hata değil. Bugünkü (\texttt{7.0.0} gibi
modern) bir kernelde yeni kod yazarken \texttt{struct mutex} tercih
edilir; kavramsal olarak burada anlatılan semaphore-mutex mantığıyla
birebir aynı işi görür, sadece daha hafif ve modern bir API'dir.
\end{notkutu}

\subsection*{B\"ol\"um 5 -- \"On Plana \c{C}\i kan Uyar\i lar}
\begin{enumerate}[leftmargin=1.2cm]
  \item Kilitleri, korunan nesneyi sistemin geri kalanına açmadan
  \textbf{önce} ilklendir.
  \item Her \texttt{down}/\texttt{spin\_lock} için, \textbf{hata
  yollarını dahil} tam olarak bir eşleşen \texttt{up}/\texttt{spin\_unlock}.
  \item Spinlock tutarken asla uyuma -- \texttt{copy\_from\_user} ve
  \texttt{kmalloc(GFP\_KERNEL)} gibi gizli uyku noktalarına dikkat.
  \item Aynı-CPU'da kesme kaynaklı kendine-kilitlenmeyi önlemek için
  gerektiğinde \texttt{\_irqsave}/\texttt{\_irq} varyantlarını kullan.
  \item Birden fazla kilidi hep \textbf{aynı sırada} al; semaforları
  spinlock'lardan önce al.
  \item BKL'yi (\texttt{lock\_kernel}) yeni kodda asla kullanma.
  \item Erken optimizasyona direnç göster -- kaba taneli kilitle
  başla, gerçek çekişme kanıtlanınca inceltmeyi düşün.
\end{enumerate}

\newpage

%=================================================================
\section{Geli\c{s}mi\c{s} Karakter S\"ur\"uc\"u \.{I}\c{s}lemleri}
\label{sec:bolum6}
%=================================================================

\begin{notkutu}{Kernel sürümü notu}
Kitap bu bölümde \textbf{eski tip} \texttt{.ioctl} imzasını kullanır:
\texttt{int (*ioctl)(struct inode *, struct file *, unsigned int,
unsigned long)}. Daha yeni kernellerde (3.x ve sonrası) bunun yerine
\texttt{.unlocked\_ioctl} kullanılır (\texttt{struct inode}
parametresi olmadan, ve örtük "Big Kernel Lock" kaldırılmış halde).
Kavramlar (komut kodlama, \texttt{access\_ok}, capabilities) birebir
aynı kalır; sadece fonksiyon imzası değişmiştir -- bunu bilerek
oku.
\end{notkutu}

\subsection{ioctl: Genel Bak\i\c{s}}

Kullanıcı tarafı imza: \texttt{int ioctl(int fd, unsigned long cmd,
...);}. Buradaki \texttt{...} gerçek bir değişken-argüman listesi
\textbf{değildir} (sistem çağrıları değişken argüman sayısı alamaz) --
sadece derleyicinin üçüncü argüman üzerinde tip kontrolünü
\textbf{kapatmak} için oradadır. Gerçekte tek bir \texttt{unsigned
long} argüman iletilir; kullanıcı bir işaretçi geçirirse kernel onu
\texttt{unsigned long} olarak alır.

Çoğu ioctl uygulaması tek bir büyük \texttt{switch(cmd)} bloğudur.

\subsection{Komut Numaras\i\ Kodlama\c{s}\i : Neden Bu Kadar Karma\c{s}\i k?}

\textbf{Temel gerekçe} (kitaptan): ioctl komut numaraları \textbf{sistem
genelinde benzersiz} olmalıdır -- aksi halde bir program yanlışlıkla
\emph{yanlış cihaza} doğru numaralı ama farklı anlamdaki bir komut
gönderebilir (örneğin bir FIFO'ya, bir seri portun baud rate komutunu
göndermek gibi). Her numara benzersizse, yanlış eşleşme
\texttt{-EINVAL} ile başarısız olur -- yanlışlıkla \textbf{istenmeyen
bir işlemi çalıştırmak yerine}.

Modern kodlama şeması dört bit alanından oluşur
(\texttt{<asm/ioctl.h>}):
\begin{itemize}[leftmargin=1.2cm]
  \item \textbf{type} (magic number) -- sürücü başına bir sayı, 8 bit.
  \item \textbf{number} -- sıra numarası, 8 bit, kendi başına bir
  anlamı yok.
  \item \textbf{direction} -- \textbf{uygulamanın bakış açısından}
  veri yönü: \texttt{\_IOC\_NONE}, \texttt{\_IOC\_READ} (uygulama
  cihazdan okuyor), \texttt{\_IOC\_WRITE} (uygulama cihaza yazıyor),
  ya da ikisinin OR'u.
  \item \textbf{size} -- veri boyutu; kernel bunu \textbf{zorunlu
  kılmaz} ama iyi pratik olarak kullanıcı hatalarını yakalamaya yarar.
\end{itemize}

Kodlama makroları:
\begin{lstlisting}
_IO(type, nr)              /* argumansiz komut */
_IOR(type, nr, datatype)   /* driver'dan OKUMA */
_IOW(type, nr, datatype)   /* driver'a YAZMA */
_IOWR(type, nr, datatype)  /* iki yonlu */
\end{lstlisting}

scull'un tanımladığı komutlar:
\begin{lstlisting}
#define SCULL_IOC_MAGIC  'k'
#define SCULL_IOCRESET    _IO(SCULL_IOC_MAGIC, 0)
#define SCULL_IOCSQUANTUM _IOW(SCULL_IOC_MAGIC,  1, int)  /* Set: pointer ile */
#define SCULL_IOCTQUANTUM _IO(SCULL_IOC_MAGIC,   3)       /* Tell: deger ile */
#define SCULL_IOCGQUANTUM _IOR(SCULL_IOC_MAGIC,  5, int)  /* Get: pointer ile */
#define SCULL_IOCQQUANTUM _IO(SCULL_IOC_MAGIC,   7)       /* Query: return degeriyle */
#define SCULL_IOC_MAXNR 14
\end{lstlisting}

\textbf{Dönüş değeri kuralı}: bilinmeyen \texttt{cmd} için POSIX'e
göre \texttt{-ENOTTY} döndürülmeli ("Inappropriate ioctl for device")
-- tarihsel olarak sadece tty'lerin ioctl'i olduğu dönemden kalma
(\texttt{-EINVAL} de pratikte sık görülür ama standart olan
\texttt{-ENOTTY}'dir).

\begin{notkutu}{Önceden tanımlı komutlar sürücünden önce yorumlanır}
Bazı ioctl komutları (\texttt{FIONBIO}, \texttt{FIOASYNC}...) kernel
tarafından \textbf{senin \texttt{.ioctl} metodun çağrılmadan önce}
yorumlanır. Eğer kendi komutun için aynı sayısal değeri seçersen,
komutun sürücüne \textbf{asla ulaşmaz} -- kernel onu yakalar, uygulama
beklenmedik davranış görür.
\end{notkutu}

\subsection{ioctl Argüman\i n\i n G\"uvenli Kullan\i m\i}

\texttt{arg} bir işaretçiyse, dereference etmeden önce doğrulanmalıdır:

\begin{lstlisting}
int access_ok(int type, const void *addr, unsigned long size);
/* type: VERIFY_READ veya VERIFY_WRITE (WRITE, READ'in ustkumesidir) */
/* DONUS: 1 = basari, 0 = basarisizlik -- coğu kernel fonksiyonunun TERSI! */
\end{lstlisting}

scull'un komut doğrulaması, hem sayısal geçerliliği hem yön bitlerini
kontrol eder:
\begin{lstlisting}
if (_IOC_TYPE(cmd) != SCULL_IOC_MAGIC) return -ENOTTY;
if (_IOC_NR(cmd) > SCULL_IOC_MAXNR) return -ENOTTY;

if (_IOC_DIR(cmd) & _IOC_READ)
    err = !access_ok(VERIFY_WRITE, (void __user *)arg, _IOC_SIZE(cmd));
else if (_IOC_DIR(cmd) & _IOC_WRITE)
    err = !access_ok(VERIFY_READ, (void __user *)arg, _IOC_SIZE(cmd));
if (err) return -EFAULT;
\end{lstlisting}

\begin{notkutu}{Kafa karıştırıcı ama önemli: yön tersine döner}
\textbf{ioctl komutunun} yönü uygulama-odaklıdır (\texttt{\_IOC\_READ}
= uygulama okuyor); \textbf{access\_ok'un} \texttt{type}'ı kernel-
odaklıdır (kernel bu belleğe \textbf{yazacaksa} \texttt{VERIFY\_WRITE}
gerekir). Uygulamanın "okuması," kernelin oraya "yazması" demektir --
bu tersine dönüş kolayca kafa karıştırır, dikkatli ol.
\end{notkutu}

Tek bir değer aktarmak için, \texttt{access\_ok} zaten geçtiyse daha
hızlı makrolar kullanılabilir:
\begin{lstlisting}
__put_user(datum, ptr);  /* access_ok'u ATLAR, onceden dogrulanmis olmali */
__get_user(local, ptr);
put_user(datum, ptr);    /* kendi icinde access_ok cagirir */
get_user(local, ptr);
\end{lstlisting}

\subsection{Capabilities: root Yerine \.{I}nce Taneli Yetki}

Geleneksel Unix modeli hep-ya-da-hiç'tir (sadece root). Linux bunun
yerine \textbf{capabilities} sunar -- ayrıcalıklı işlemler alt
gruplara bölünür. Sürücü yazarları için önemli olanlar:
\texttt{CAP\_SYS\_ADMIN} (genel yönetim, "daha spesifik bir capability
yokluğunda" en sık kullanılan), \texttt{CAP\_SYS\_MODULE},
\texttt{CAP\_SYS\_RAWIO}, \texttt{CAP\_DAC\_OVERRIDE}.

\begin{lstlisting}
int capable(int capability);  /* surecin bu yetkiye sahip olup olmadigini dondurur */

case SCULL_IOCSQUANTUM:
    if (! capable (CAP_SYS_ADMIN))
        return -EPERM;
    retval = __get_user(scull_quantum, (int __user *)arg);
    break;
\end{lstlisting}

scull'un politikası: herkes quantum/qset boyutunu \textbf{sorgulayabilir}
(Get/Query), ama sadece yetkili kullanıcılar \textbf{değiştirebilir}
(Set/Tell) -- kötü değerler performansı bozabilir.

\subsection{Blocking I/O: Process'i Uyutmak}

\subsubsection{Üç temel kural (kitaptan, aynen)}

\begin{enumerate}[leftmargin=1.2cm]
  \item \textbf{Atomik bağlamda asla uyuma.} Bir spinlock, seqlock,
  ya da RCU kilidi tutarken, ya da kesmeler kapalıyken uyuyamazsın.
  Bir semafor tutarken uyumak \textbf{yasal}dır ama çok kısa olmalı ve
  o semaforu tutmanın seni uyandıracak tarafı bloklamadığından emin
  olmalısın.
  \item \textbf{Uyandıktan sonra sistem durumu hakkında hiçbir varsayım
  yapma.} Ne kadar uyuduğunu, o sırada neyin değiştiğini bilemezsin;
  başka bir process aynı olay için uyuyor olabilir ve senden önce
  kaynağı "kazanmış" olabilir. \textbf{Beklediğin koşulu her zaman
  yeniden kontrol etmelisin.}
  \item \textbf{Process'in uyanacağından emin olmadan uyuma.} Bir
  yerde, birinin seni uyandırması garanti olmalı.
\end{enumerate}

\subsubsection{wait\_event ailesi}

\begin{lstlisting}
DECLARE_WAIT_QUEUE_HEAD(name);          /* derleme-zamani */
wait_queue_head_t q; init_waitqueue_head(&q);  /* calisma-zamani */

wait_event(queue, condition);
wait_event_interruptible(queue, condition);  /* tercih edilen */
wait_event_timeout(queue, condition, zamanasimi);
wait_event_interruptible_timeout(queue, condition, zamanasimi);
\end{lstlisting}

\begin{notkutu}{wait\_event neden race-condition'a dayanıklı}
\texttt{condition}, makro tarafından \textbf{hem uyumadan önce hem
uyandıktan sonra} değerlendirilir ve koşul doğru olana kadar makro
içeride \textbf{döngüyle} beklemeye devam eder. Bu, "check sonra
uyu" arasındaki klasik race'i (kaçırılan uyandırma) kapatan yerleşik
mekanizmadır -- ayrıca \texttt{condition}'ın \textbf{yan etkisi
olmamalı}, çünkü kaç kez değerlendirileceği belirsizdir.
\texttt{wait\_event\_interruptible} sinyalle kesilirse sıfırdan farklı
döner -- bu durumda sürücü \texttt{-ERESTARTSYS} döndürmelidir.
\end{notkutu}

\begin{lstlisting}
void wake_up(wait_queue_head_t *queue);
void wake_up_interruptible(wait_queue_head_t *queue);
\end{lstlisting}

\subsubsection{sleepy örneği: paylaşılan bayrak race'i}

\begin{lstlisting}
static DECLARE_WAIT_QUEUE_HEAD(wq);
static int flag = 0;

ssize_t sleepy_read (struct file *filp, char __user *buf, size_t count, loff_t *pos)
{
    wait_event_interruptible(wq, flag != 0);
    flag = 0;
    return 0;
}
ssize_t sleepy_write (struct file *filp, const char __user *buf, size_t count, loff_t *pos)
{
    flag = 1;
    wake_up_interruptible(&wq);
    return count;
}
\end{lstlisting}

\begin{notkutu}{Kitabın açıkça gösterdiği gotcha}
İki process \texttt{sleepy\_read}'de uyuyorken \texttt{sleepy\_write}
bir kez çalışırsa, \texttt{wake\_up\_interruptible} \textbf{ikisini de}
uyandırır -- ikisi de \texttt{flag != 0}'ı sıfırlanmadan önce görebilir,
yani ikisi de yanlışlıkla "koşul doğru" sonucuna varabilir. Kitap:
"gerçek bir sürücüde bu tür bir race, teşhisi zor, nadir çökmeler
yaratabilir." Tam olarak bu yüzden \texttt{scullpipe} bir bayrak
yerine \textbf{semafor + dairesel tampon} kullanır -- durumu atomik
olarak test eder.
\end{notkutu}

\subsubsection{O\_NONBLOCK ve -EAGAIN}

\texttt{filp->f\_flags} içindeki \texttt{O\_NONBLOCK} bayrağı
ayarlıysa, veri/yer yokken \texttt{read}/\texttt{write} \textbf{hemen}
\texttt{-EAGAIN} döndürmelidir -- process'i asla bloklamaz. Sadece
\texttt{read}, \texttt{write}, ve \texttt{open} bu bayraktan etkilenir.

\subsubsection{scullpipe: Ger\c{c}ek Bir Bloklayan Okuma \"Ornek}

\begin{lstlisting}
struct scull_pipe {
    wait_queue_head_t inq, outq;  /* okuma ve yazma kuyruklari */
    char *buffer, *end;
    int buffersize;
    char *rp, *wp;                /* nereden oku, nereye yaz */
    int nreaders, nwriters;
    struct semaphore sem;
    struct cdev cdev;
};
\end{lstlisting}

\textbf{İki ayrı wait queue}: \texttt{inq} (okuyucular veri bekler),
\texttt{outq} (yazıcılar yer bekler) -- kuralda geçen "farklı işlemler
için farklı kuyruk" ilkesinin somut hali.

\begin{lstlisting}
static ssize_t scull_p_read (struct file *filp, char __user *buf,
                              size_t count, loff_t *f_pos)
{
    struct scull_pipe *dev = filp->private_data;
    if (down_interruptible(&dev->sem))
        return -ERESTARTSYS;

    while (dev->rp == dev->wp) { /* okunacak bir sey yok */
        up(&dev->sem);           /* KILIDI BIRAK -- yoksa yazicilar bizi hic uyandiramaz */
        if (filp->f_flags & O_NONBLOCK)
            return -EAGAIN;
        if (wait_event_interruptible(dev->inq, (dev->rp != dev->wp)))
            return -ERESTARTSYS;
        if (down_interruptible(&dev->sem))  /* uyandik: kilidi TEKRAR al */
            return -ERESTARTSYS;
        /* dongu basina donup kosulu YENIDEN kontrol et */
    }
    /* ...veriyi kopyala, rp'yi ilerlet... */
    up(&dev->sem);
    wake_up_interruptible(&dev->outq);  /* yazicilari uyandir */
    return count;
}
\end{lstlisting}

\begin{notkutu}{Bu kodun öğrettiği en önemli desen}
Tampon boşsa, semafor \textbf{uyumadan önce bırakılır} -- kitabın
kendi ifadesiyle: "eğer onu tutarken uyusaydık, hiçbir yazıcının bizi
uyandırma fırsatı olmazdı" (doğrudan self-deadlock). Uyandıktan sonra:
sinyal geldiyse \texttt{-ERESTARTSYS}; değilse semafor
\textbf{yeniden alınır} ve \texttt{while} döngüsü koşulu \textbf{yeniden
test eder} -- çünkü uyanman veri orada olduğu anlamına gelmez, başka
bir okuyucu senden önce veriyi almış olabilir. Bu, Kural 2'nin (asla
uyanma sonrası varsayım yapma) birebir kod haline gelmiş hali.
\end{notkutu}

\subsubsection{Elle Uyuma Mekaniği (wait\_event'in altında ne oluyor)}

\begin{enumerate}[leftmargin=1.2cm]
  \item Bir \texttt{wait\_queue\_t} girdisi oluşturulup kuyruğa eklenir.
  \item Process durumu değiştirilir: \texttt{TASK\_RUNNING} (çalışabilir),
  \texttt{TASK\_INTERRUPTIBLE}/\texttt{TASK\_UNINTERRUPTIBLE}
  (uyuyor). \textbf{Sadece durumu değiştirmek process'i uyutmaz} --
  zamanlayıcıya nasıl davranacağını söyler, henüz CPU bırakılmamıştır.
  \item \textbf{Kritik desen}: durum ayarlandıktan hemen sonra, ama
  \texttt{schedule()} çağrılmadan önce koşul tekrar kontrol edilir:
  \begin{lstlisting}
if (!condition)
    schedule();
  \end{lstlisting}
  Bu, olayın \emph{ne zaman} gerçekleştiğine bakılmaksızın doğru
  çalışır: koşul durum ayarlanmadan \textbf{önce} doğru olduysa, bu
  kontrolde yakalanır (uyumadan atlanır); \textbf{sonra} gerçekleşirse,
  process zaten çalıştırılabilir yapılmıştır.
  \item \texttt{schedule()} -- zamanlayıcıyı çağırır, CPU'yu bırakır.
\end{enumerate}

\begin{notkutu}{Tarihi bir uyarı: sleep\_on kullanma}
\texttt{sleep\_on()}/\texttt{interruptible\_sleep\_on()} kesinlikle
\textbf{kullanımdan kaldırılmıştır} (deprecated) -- race condition'a
karşı \textbf{hiçbir koruma sağlamazlar}: koşulu kontrol etmekle
gerçekten uyumak arasında her zaman bir pencere vardır, ve o pencerede
gelen bir uyandırma kaybolur. Kitap: "\texttt{sleep\_on} çağıran kod
asla tam olarak güvenli değildir."
\end{notkutu}

\subsection{poll ve select}

Bir process'in bloklamadan "okuyabilir miyim / yazabilir miyim"
sorabilmesini sağlayan mekanizma. Tek bir sürücü metoduyla desteklenir:

\begin{lstlisting}
unsigned int (*poll) (struct file *filp, poll_table *wait);
\end{lstlisting}

İki görevi vardır: (1) \texttt{poll\_wait()}'i durumdaki değişikliği
gösterebilecek her wait queue için çağırmak, (2) hangi işlemlerin
\textbf{bloklamadan} tamamlanabileceğini gösteren bir bit maskesi
döndürmek (\texttt{POLLIN|POLLRDNORM} okunabilir,
\texttt{POLLOUT|POLLWRNORM} yazılabilir, \texttt{POLLHUP} EOF,
\texttt{POLLERR} hata).

\begin{lstlisting}
static unsigned int scull_p_poll(struct file *filp, poll_table *wait)
{
    struct scull_pipe *dev = filp->private_data;
    unsigned int mask = 0;

    down(&dev->sem);
    poll_wait(filp, &dev->inq,  wait);
    poll_wait(filp, &dev->outq, wait);
    if (dev->rp != dev->wp)
        mask |= POLLIN | POLLRDNORM;
    if (spacefree(dev))
        mask |= POLLOUT | POLLWRNORM;
    up(&dev->sem);
    return mask;
}
\end{lstlisting}

\begin{notkutu}{write asla veri iletimini beklemek için bloklamamalı}
Kitabın vurguladığı önemli bir kural: uygulamalar \texttt{select}'i
tam olarak "bu \texttt{write} bloklayacak mı" sorusuna cevap almak
için kullanır -- eğer \texttt{poll} yazılabilir dediyse,
\texttt{write} \textbf{gerçekten} bloklamamalıdır (\texttt{O\_NONBLOCK}
ayarlı olmasa bile). Verinin fiziksel olarak iletildiğinden emin olmak
isteyen bir uygulama bunun yerine \texttt{.fsync} kullanmalıdır.
\end{notkutu}

\subsection{Asenkron Bildirim (fasync / SIGIO)}

Kullanıcı tarafı iki adımlı bir \texttt{fcntl} anlaşmasıdır:
\begin{lstlisting}
fcntl(fd, F_SETOWN, getpid());              /* 1. kimi bilgilendirecegini soyle */
fcntl(fd, F_SETFL, fcntl(fd,F_GETFL)|FASYNC); /* 2. bildirimi ac */
\end{lstlisting}
Ardından kernel, dosyada yeni veri geldiğinde sahibe \texttt{SIGIO}
sinyali gönderir. Sürücü tarafında hazır yardımcı fonksiyonlar
kullanılır:
\begin{lstlisting}
int  fasync_helper(int fd, struct file *filp, int mode, struct fasync_struct **fa);
void kill_fasync(struct fasync_struct **fa, int sig, int band);
\end{lstlisting}
scull örneği:
\begin{lstlisting}
static int scull_p_fasync(int fd, struct file *filp, int mode)
{
    struct scull_pipe *dev = filp->private_data;
    return fasync_helper(fd, filp, mode, &dev->async_queue);
}
/* write icinde, yeni veri geldiginde: */
if (dev->async_queue)
    kill_fasync(&dev->async_queue, SIGIO, POLL_IN);
\end{lstlisting}

\begin{notkutu}{Kolayca unutulan temizlik adımı}
\texttt{.fasync} metodu, dosya \textbf{kapanırken de} çağrılmalıdır
(genelde \texttt{.release} içinde, \texttt{scull\_p\_fasync(-1, filp,
0);} gibi) -- yoksa kapanmış bir \texttt{filp} bildirim listesinde
asılı kalır.
\end{notkutu}

\subsection{Cihaz\i\ Aramaya Kapatmak: nonseekable\_open}

\texttt{.llseek} \texttt{NULL} bırakılırsa, kernelin \textbf{varsayılan
davranışı yine de aramaya (seek) izin verir} -- yani seri port gibi
gerçek bir veri \emph{akışı} sunan (veri \emph{alanı} değil) bir cihaz
için, aramayı devre dışı bırakmak istiyorsan \texttt{.llseek}'i boş
bırakmak \textbf{yetmez}. Açıkça reddetmek gerekir:

\begin{lstlisting}
int nonseekable_open(struct inode *inode, struct file *filp);
/* .open icinde cagir, ve fops'ta: */
.llseek = no_llseek,
\end{lstlisting}

scull, gerçek bir veri alanı olduğu için kendi \texttt{.llseek}'ini
uygular -- \texttt{SEEK\_END} durumunda cihazın güncel boyutunu
(\texttt{dev->size}) kullanarak yeni konumu hesaplar.

\subsection{Cihaz Dosyas\i\ \"Uzerinde Eri\c{s}im Kontrol\"u}

Dosya izin bitleri \textbf{kimin} bir cihazı açabileceğini kısıtlar --
ama bazen yetkili kullanıcılar arasında bile \textbf{aynı anda sadece
biri} açık tutabilmelidir. Kitap dört farklı politikayı ayrı scull
varyantları olarak gösterir:

\begin{itemize}[leftmargin=1.2cm]
  \item \textbf{scullsingle} -- \texttt{atomic\_t} sayaç ile
  ("brute-force," kitabın kendi tabiriyle "genelde kaçınılması
  gereken" bir yaklaşım -- meşru eşzamanlı kullanım desenlerini
  (örn. bir okuma-durumu process'i + bir yazma-verisi process'i)
  engeller):
  \begin{lstlisting}
static atomic_t scull_s_available = ATOMIC_INIT(1);
if (! atomic_dec_and_test (&scull_s_available)) {
    atomic_inc(&scull_s_available);
    return -EBUSY;
}
  \end{lstlisting}
  \item \textbf{sculluid} -- sadece \textbf{tek bir kullanıcı}
  (birden fazla process'e izin vererek) açabilir; spinlock ile
  korunan sayaç + sahip uid:
  \begin{lstlisting}
spin_lock(&scull_u_lock);
if (scull_u_count &&
        (scull_u_owner != current->uid) &&
        (scull_u_owner != current->euid) &&
        !capable(CAP_DAC_OVERRIDE)) {
    spin_unlock(&scull_u_lock);
    return -EBUSY;   /* -EPERM kullaniciyi yanlis yone yonlendirir */
}
  \end{lstlisting}
  \begin{notkutu}{-EBUSY neden -EPERM değil? (bilinçli UX kararı)}
  Kitap açıkça bunun kasıtlı olduğunu söylüyor: "kullanıcılar
  'Permission denied' mesajını dosya izin/sahiplik bitlerini kontrol
  etme daveti olarak yorumlar" -- oysa asıl sorun bu değildir.
  "Device busy" doğru şekilde "zaten kullanan bir process ara"
  fikrini verir.
  \end{notkutu}
  \item \textbf{scullwuid} -- sculluid gibi ama \texttt{-EBUSY}
  yerine bir wait queue'da \textbf{bekler}; cron gibi arka plan
  görevleri için etkileşimli redden daha uygun olabilir. Kitap
  gerçek dünya örneği verir: Linux teyp sürücüsü, aynı fiziksel
  cihaz için \textbf{farklı davranışlar} sunan birden fazla cihaz
  düğümü (device node) kullanır -- bloklama-vs-reddetme gibi uyumsuz
  politikaların her ikisine de ihtiyaç varsa, çözüm genelde tek bir
  cihazı iki politikayla zorlamak değil, \textbf{iki ayrı cihaz
  düğümü} sunmaktır.
  \item \textbf{scullpriv} -- her açan process'e (kontrol eden
  terminaline göre anahtarlanan) \textbf{ayrı bir sanal kopya}
  sunar -- sadece scull gibi saf yazılım cihazlarında mümkün. Anahtar
  olarak terminal yerine \texttt{uid} kullanılsaydı, her kullanıcıya
  özel bir sanal cihaz elde edilirdi -- bu tamamen bir politika
  seçimidir.
\end{itemize}

\subsection*{B\"ol\"um 6 -- \"On Plana \c{C}\i kan Uyar\i lar}
\begin{enumerate}[leftmargin=1.2cm]
  \item ioctl komut numaraları asla küçük ardışık tam sayılar
  olmamalı -- type/number/direction/size kodlamasına uyulmalı.
  \item \texttt{access\_ok}'un \texttt{VERIFY\_READ}/\texttt{VERIFY\_WRITE}
  anlamı, ioctl'in kendi yön bitlerine göre \textbf{tersine} döner.
  \item \texttt{\_\_put\_user}/\texttt{\_\_get\_user} sadece
  \texttt{access\_ok} ile önceden doğrulanmış adreslerde kullanılmalı.
  \item Check-then-sleep race'i: koşulu kontrol edip sonra uyumak iki
  \textbf{ayrı, senkronize olmayan} adımsa, arada gelen bir uyandırma
  sonsuza kadar kaybolabilir -- \texttt{wait\_event*} bunu içeride
  otomatik çözer, elle kod yazıyorsan sen çözmelisin.
  \item Beklemeden önce, uyandırıcı tarafın ihtiyaç duyduğu kilidi
  \textbf{bırak}.
  \item Sinyalle kesilme her zaman \texttt{-ERESTARTSYS} ile
  bildirilmeli.
  \item \texttt{.llseek}'i boş bırakmak aramayı devre dışı
  \textbf{bırakmaz} -- aramayı reddetmek için açıkça
  \texttt{nonseekable\_open} + \texttt{no\_llseek} gerekir.
\end{enumerate}

\newpage

%=================================================================
\section{Bellek Ay\i rma}
\label{sec:bolum8}
%=================================================================

Senin \texttt{ahmet.c}'in şu an \texttt{static char buffer[1024]} ile
\textbf{derleme-zamanı} sabit bir alan kullanıyor. Bu bölüm, kernel'in
\textbf{çalışma-zamanında} bellek istemenin kaç farklı yolu olduğunu ve
hangisini ne zaman seçmen gerektiğini anlatır.

\subsection{kmalloc: Kernel'in malloc'u, Ama Farklı Kurallarla}

\texttt{kmalloc(size\_t size, int flags)} çağrısı \textbf{fiziksel
olarak bitişik} (contiguous) bellek döndürür -- ve \textbf{önceki
içeriğini temizlemez}; user space'e gönderilecek ya da cihaza
yazılacaksa elle sıfırlanmalı (aksi halde önceki kernel verisi sızabilir).

En önemli iki bayrak, ne zaman hangisinin kullanılacağı konusunda net
bir kural taşır:

\begin{itemize}[leftmargin=1.2cm]
  \item \textbf{\texttt{GFP\_KERNEL}} -- normal, process bağlamındaki
  (process context) çağrılar için. Bellek azsa \textbf{uyuyabilir}
  (process'i bekletir). Bu yüzden onu çağıran fonksiyon
  \textbf{reentrant olmalı ve atomik bağlamda çalışmamalıdır.}
  \item \textbf{\texttt{GFP\_ATOMIC}} -- process bağlamı
  \textbf{dışında} (kesme işleyicisi, tasklet, timer) zorunludur.
  \textbf{Asla uyumaz}; kernel bunun için ayrı bir acil-durum
  rezervi tutar, o da tükenirse ayırma \textbf{başarısız olur}
  (ama hiç uyumaz).
\end{itemize}

\begin{notkutu}{En kritik kural}
Kesme bağlamında (ya da Bölüm 5'te öğrendiğin bir spinlock tutarken)
\texttt{GFP\_KERNEL} kullanmak ciddi bir hatadır -- kod uyuyabilir,
ve zaten uyuyamayacağın bir yerdesin. Doğru bayrak
\texttt{GFP\_ATOMIC}'tir.
\end{notkutu}

\textbf{Boyut sınırı}: kmalloc, önceden tanımlı sabit boyut sınıflarından
(2'nin katları) bellek verir -- istediğinden biraz \textbf{fazlasını}
alabilirsin (en kötü ihtimalle 2 katı). Taşınabilir kod
\textbf{128 KB}'tan fazlasını kmalloc ile isteyemez -- sistem
çalıştıkça bitişik büyük blok bulmak zorlaşır (parçalanma/fragmentation).

\subsection{Slab Cache: Ayn\i\ Boyuttaki Nesneleri S\i k S\i k Ay\i r\i p B\i rakanlar \.{I}\c{c}in}

Bir sürücü \textbf{aynı boyutta} nesneleri sürekli ayırıp serbest
bırakıyorsa (USB ve SCSI sürücüleri gibi), kmalloc yerine
\texttt{kmem\_cache\_create()} ile kendi havuzunu (lookaside cache)
oluşturabilir -- \texttt{kmem\_cache\_alloc}/\texttt{kmem\_cache\_free}
ile ayırıp bırakır, \texttt{kmem\_cache\_destroy} ile kapatır (sadece
tüm nesneler geri döndüyse başarılı olur -- başarısız dönüş, bir
sızıntı olduğunun işaretidir). Kitabın \texttt{scullc} varyantı bunu
gösterir: kmalloc yerine cache kullanmak hız değil, \textbf{daha
öngörülebilir/verimli bellek kullanımı} sağlar.

\subsection{Sayfa Bazl\i\ Ayirma ve vmalloc: Farkl\i\ İhtiyaçlar \.{I}\c{c}in Farkl\i\ Ara\c{c}lar}

\texttt{\_\_get\_free\_pages()} ailesi, kmalloc'un aksine tam sayfa(lar)
ister -- büyük parçalar için kmalloc'tan daha verimlidir (kmalloc'un
öngörülemeyen fazlalığı yoktur).

\begin{notkutu}{vmalloc ile kmalloc arasındaki en önemli fark}
\texttt{vmalloc()}'in döndürdüğü bellek \textbf{sanal adres uzayında
bitişiktir} ama \textbf{fiziksel bellekte bitişik olmak zorunda
değildir} (her sayfa ayrı ayrı alınıp sanal olarak yan yana
haritalanır). \texttt{kmalloc} ikisini de garanti eder. \textbf{Donanım
DMA yapacaksa \texttt{vmalloc} kullanılamaz} -- donanım fiziksel
adrese ihtiyaç duyar, vmalloc'un sanal adresi donanım için anlamsızdır.
Kitap açıkça caydırıyor: "çoğu durumda vmalloc kullanımı önerilmez."
Uygun olduğu nadir durum: yalnızca yazılımda var olan, büyük, ardışık
bir tampon (örn. modül yükleme mekanizmasının kendisi, modül kodunu
yüklerken vmalloc kullanır).
\end{notkutu}

\subsection{CPU-Ba\c{s}\i na De\u{g}i\c{s}kenler: Kilitsiz H\i z}

2.6 kernel'in bir özelliği: \texttt{DEFINE\_PER\_CPU(tip, isim)} ile
tanımlanan bir değişkenin \textbf{her işlemci kendi kopyasına} sahip
olur -- kilitleme gerekmeden hızlı sayaç/istatistik tutmak için
idealdir (ağ paket sayaçları gibi).

\begin{notkutu}{get\_cpu\_var / put\_cpu\_var eşleşmesi zorunlu}
\texttt{get\_cpu\_var(isim)} preemption'ı (önceliklenmeyi) devre dışı
bırakır ve o CPU'nun kopyasına erişim verir; \texttt{put\_cpu\_var(isim)}
\textbf{bunu tekrar etkinleştirir}. İkisi her zaman eşleşmelidir --
\texttt{put\_cpu\_var}'ı atlamak preemption'ı kalıcı olarak kapalı
bırakır.
\end{notkutu}

\subsection{Büyük Tamponlar: Neden Kaçınmalı}

Kitap açık bir tavsiye veriyor: büyük, bitişik bellek ayırmaları
\textbf{zamanla arızaya eğilimlidir} (sistem çalıştıkça bellek
parçalanır). Gerçekten büyük bir tampon gerekiyorsa, en iyi çözüm
genelde tek bir dev blok istemek yerine \textbf{scatter/gather}
işlemleridir (Bölüm 15'in konusu, bu belgenin kapsamı dışında).

\begin{ahmetkutu}
Bu bölüm, README'deki "sırada" listesindeki maddeyi netleştiriyor:
\texttt{ahmet\_init} içinde \texttt{buffer = kmalloc(buffer\_size,
GFP\_KERNEL);} ile dinamik ayırıp \texttt{ahmet\_exit}'te
\texttt{kfree(buffer);} ile serbest bırakmak -- 1024 byte gibi küçük
bir tampon için pratikte fark yaratmaz, ama \textbf{doğru API'yi doğru
yerde kullanmayı} öğrenmiş olursun: sürücün process bağlamında
(\texttt{ahmet\_init}, bir modül yükleme işlemidir) çalıştığı için
burada \texttt{GFP\_KERNEL} doğru seçimdir.
\end{ahmetkutu}

\newpage

%=================================================================
\section{Kesme (Interrupt) Y\"onetimi}
\label{sec:bolum10}
%=================================================================

\begin{notkutu}{Neden bu bölüm senin için hemen uygulanabilir değil, ama önemli}
\texttt{ahmet.c} gerçek donanıma bağlı değil -- kesme (interrupt)
üretmiyor. Ama gömülü sistem stajında er ya da geç gerçek bir
donanımla (bir buton, bir sensör, bir seri port çipi) çalışacaksın,
ve o donanım seni \textbf{kesme ile} uyaracak. Bu bölümü şimdi
kavramsal olarak anlamak, o günü daha rahat karşılamanı sağlar.
\end{notkutu}

\subsection{Kesme Neden Var: Polling'e Alternatif}

Donanım, "bir şey oldu" demek için CPU'yu sürekli yokladırmak
(polling -- boşa CPU harcar) yerine bir \textbf{kesme sinyali}
gönderir; CPU o an ne yapıyorsa bırakıp senin kayıtlı fonksiyonunu
çalıştırır. Kesme hatları (IRQ) sınırlı bir kaynaktır -- kullanmadan
önce talep edilmeli:

\begin{lstlisting}
int request_irq(unsigned int irq, irq_handler_t handler,
                 unsigned long flags, const char *dev_name, void *dev_id);
void free_irq(unsigned int irq, void *dev_id);
\end{lstlisting}

\begin{notkutu}{Kernel sürümü uyarısı: bu bölüm de eski API kullanıyor}
Kitaptaki \texttt{handler} imzası \texttt{struct pt\_regs *regs} adlı
eski, artık kaldırılmış bir üçüncü parametre içerir -- modern
kernellerde yok. Bayraklar da \texttt{SA\_INTERRUPT}/\texttt{SA\_SHIRQ}
olarak geçer; bunlar modern kernellerde \texttt{IRQF\_SHARED} gibi
isimlere dönüştü (\texttt{SA\_INTERRUPT}'ın karşılığı olan "fast
handler" kavramı tamamen kaldırıldı). Kavramlar aynı, sadece isimler
değişti.
\end{notkutu}

\texttt{dev\_id}'nin iki görevi vardır: (1) \texttt{free\_irq}'nun
\textbf{hangi} kaydı sileceğini bilmesi, (2) \textbf{paylaşılan}
bir IRQ hattında, kesme geldiğinde hangi cihaz örneğinin ilgilendiğini
ayırt etmek. Bir handler'ın dönüş değeri: \texttt{IRQ\_HANDLED}
(cihazım gerçekten kesme üretti, ilgilendim) ya da \texttt{IRQ\_NONE}
(bu kesme benim cihazımdan gelmedi).

\subsection{Handler Kurallar\i : Neden Bu Kadar K\i s\i tl\i}

Bir kesme işleyicisi, Bölüm 5'teki spinlock kurallarının en katı
haliyle çalışır: \textbf{user space'e erişemez, uyuyamaz,
\texttt{schedule()} çağıramaz}. Bu yüzden handler'ın işi minimal
tutulmalı -- donanımdan gelen veriyi bir tampona kaydetmek, gerekiyorsa
donanıma "aldım" demek, ve ağır işi \textbf{ertelemek}.

\begin{notkutu}{Top half / bottom half -- Bölüm 7'nin burada karşılığı}
Kayıtlı handler'ın kendisi \textbf{top half}'tir -- hızlı olmalı.
Asıl (uzun sürebilecek) iş, Bölüm 7'de gördüğün \textbf{tasklet}
ya da \textbf{workqueue} ile \textbf{bottom half}'e ertelenir. Tipik
akış: top half veriyi kapıp bir bayrak/kuyruk günceller ve
\texttt{tasklet\_schedule()}/\texttt{schedule\_work()} çağırıp
hemen döner; bottom half asıl işi (biçimlendirme, kopyalama) daha
"güvenli" bir zamanda, kesmeler açıkken yapar.
\end{notkutu}

\subsection{Kesme ile Bloklayan I/O'nun Bulu\c{s}tu\u{g}u Yer}

Bölüm 6'da öğrendiğin bloklayan \texttt{read}'in gerçek dünyadaki
karşılığı tam olarak budur: veri yokken process
\texttt{wait\_event\_interruptible} ile uyur; veri geldiğinde
\textbf{kesme işleyicisi} minimum işi yapar ve
\texttt{wake\_up\_interruptible()}'ı çağırarak bekleyen process'i
uyandırır. Yani Bölüm 6'daki "kim uyandıracak?" sorusunun donanım
dünyasındaki cevabı: \textbf{kesme işleyicisi}.

\subsection{Payla\c{s}\i lan Kesmeler}

Birden fazla cihaz aynı IRQ hattını paylaşabilir
(\texttt{IRQF\_SHARED}/\texttt{SA\_SHIRQ}). Kesme geldiğinde \textbf{o
hatta kayıtlı tüm handler'lar çağrılır} -- bu yüzden her handler önce
\textbf{kendi cihazının} gerçekten kesme ürettiğini (bir durum
register'ını okuyarak) kontrol etmeli, değilse hemen
\texttt{IRQ\_NONE} dönmelidir. \texttt{dev\_id} bu yüzden paylaşımlı
hatlarda \textbf{asla NULL olamaz} -- kernel onu handler'ları
birbirinden ayırt etmek için kullanır.

\begin{notkutu}{Kayıp kesmeler gerçek bir donanım vakasıdır}
Kitabın açık ifadesi: "Donanımla uğraşmanın gerçeği şu ki, ara sıra
cihazdan bir kesmeyi kaçırabilirsin. Bu olduğunda, sürücünün sistem
yeniden başlatılana kadar sonsuza dek durmasını \textbf{istemezsin}."
Bu yüzden gerçek sürücüler genelde bir "kaçırılan kesme" emniyet
zamanlayıcısı (watchdog timer) kurar -- zamanlayıcı, beklenen kesme
gelmediğinde handler'ı elle tetikler.
\end{notkutu}

\subsection*{B\"ol\"um 10 -- \"On Plana \c{C}\i kan Uyar\i lar}
\begin{itemize}[leftmargin=1.2cm]
  \item Modül kaldırılırken \texttt{free\_irq()} çağrılmazsa, bir
  sonraki kesme artık geçersiz/kaldırılmış koda işaret eder -- klasik
  bir use-after-free türü hata.
  \item Handler'da fazla iş yapmak, sadece o cihazı değil, sistemin
  genel tepki süresini bozar -- ağır iş her zaman ertelenmeli.
  \item Paylaşımlı bir hatta \texttt{enable\_irq}/\texttt{disable\_irq}
  çağırmak \textbf{diğer cihazları} da etkiler -- yapılmamalı.
  \item Handler ile process-bağlamındaki kod aynı veriyi paylaşıyorsa,
  process tarafı \texttt{spin\_lock\_irqsave} kullanmalı (Bölüm 5) --
  aksi halde aynı CPU'da kendini kilitleyebilirsin.
\end{itemize}

\newpage

%=================================================================
\section{The Linux Device Model (Ayg\i t Modeli)}
\label{sec:bolum14}
%=================================================================

\begin{notkutu}{Bu bölüm neden en doğrudan sana bağlı olan bölüm}
Sen zaten \texttt{class\_create()} ve \texttt{device\_create()}
kullanıyorsun -- \texttt{/dev/ahmet}'in otomatik oluşmasını sağlayan
mekanizma. Bu bölüm o "sihrin" perde arkasında ne olduğunu açıklıyor.
\end{notkutu}

\subsection{kobject: Her \c{S}eyin Alt\i nda Yatan Ortak \.{I}skelet}

Kernel, cihaz modelindeki hemen her nesneyi (\texttt{struct cdev} dahil
-- Bölüm 3'te gördüğün!) bir \texttt{struct kobject} \textbf{gömerek}
inşa eder. \texttt{kobject}'in tek işi üç şeydir: \textbf{referans
sayımı} (kaç yerin bu nesneye ihtiyacı var), \textbf{sysfs'te bir
dizin olarak görünme}, ve \textbf{hotplug olayı üretme}.

\begin{notkutu}{En kritik kural: bir kobject içeren struct'ı asla doğrudan kfree etme}
Bir kobject'in referans sayısı, senin kontrolünde olmayan sebeplerle
(örn. bir user-space process'in sysfs'teki bir dosyayı açık tutması)
sıfırdan büyük kalabilir. Bu yüzden kernel, "bu nesneyi sen değil,
\textbf{ben (kernel) son referans bittiğinde} sana haber vereceğim"
der -- bu haber verme, \texttt{kobj\_type}'a bağlı bir \texttt{release}
callback'idir. \textbf{Her kobject içeren struct'ın bir release
fonksiyonu olmak zorundadır}; \texttt{struct device} için bu yoksa
kernel "korkutucu şikayetler" basar. Gerçek bellek serbest bırakma
(\texttt{kfree}) sadece bu callback'in içinde yapılmalı --
\texttt{kobject\_put()} referans sayısını sıfıra indirdiğinde otomatik
tetiklenir.
\end{notkutu}

\subsection{sysfs: Neden Her kobject Bir Dizindir}

\texttt{kobject}'in sistemdeki konumu (\texttt{parent} işaretçisi
üzerinden) doğrudan \texttt{/sys} altında bir \textbf{dizine} eşlenir.
Bir kobject'in sunduğu her \textbf{attribute} (özellik), o dizinde bir
\textbf{dosya} olarak görünür -- \texttt{show}/\texttt{store}
fonksiyonlarıyla okunur/yazılır. Kitabın vurguladığı tasarım kuralı:
\textbf{"her dosya tek bir, insan tarafından okunabilir değer
içermeli"} -- büyük/karmaşık veri birden fazla dosyaya bölünmeli.

\subsection{Hotplug Olay\i : udev'in /dev D\"ug\"um\"unu Nas\i l Olu\c{s}turdu\u{g}u}

Bir kobject sisteme eklendiğinde/çıkarıldığında kernel bir
\textbf{hotplug olayı} (uevent) üretir -- bu, user-space'teki
\textbf{udev}'i tetikler. \textbf{Mekanizmanın tam olarak senin
\texttt{/dev/ahmet}'inle ilgisi burada}: udev, bildirim aldığında
\texttt{/sys/class/...} ağacında \texttt{dev} adlı bir dosya arar --
bu dosya major:minor numarasını içerir. \textbf{Bir sürücünün udev
için yapması gereken tek şey, kontrol ettiği cihazların major/minor
numaralarını sysfs üzerinden dışa açmaktır} -- \texttt{class\_create}
+ \texttt{device\_create} çağrıların tam olarak bunu yapıyor: bir
sınıf (\texttt{/sys/class/ahmet\_class}) oluşturup, altına \texttt{dev}
dosyasını içeren bir cihaz girdisi ekliyor.

\subsection{Bus $\to$ Device $\to$ Driver: PCI/USB'de G\"ord\"u\u{g}\"un\"un Genelle\c{s}mi\c{s} Hali}

Bir \textbf{bus} (\texttt{struct bus\_type}), üzerindeki cihazları ve
sürücüleri iki kset olarak tutar; bir \texttt{match()} fonksiyonu
hangi sürücünün hangi cihaza uyduğuna karar verir. Eşleşme bulunduğunda
sürücünün \texttt{probe()} fonksiyonu çağrılır -- Bölüm 12/13'te PCI ve
USB için gördüğün "cihaz tak, doğru sürücü otomatik eşleşsin" deseninin
\textbf{soyut, genel çerçevesi} budur.

\subsection{class: class\_create'in ge\c{c}mi\c{s}i ve senin kodunla fark\i}

\begin{ahmetkutu}
Bu bölümdeki en doğrudan sana ait bulgu: \textbf{kitabın bu baskısı
\texttt{class\_create()} diye bir fonksiyon hiç kullanmıyor} --
onun yerine (2.6.10 döneminde) \texttt{class\_simple\_create(struct
module *owner, char *name)} ve \texttt{class\_simple\_device\_add(...)}
kullanıyor. API zaman içinde şöyle evrildi:
\begin{center}
\texttt{class\_simple\_create(owner, name)} $\to$
\texttt{class\_create(owner, name)} $\to$
\texttt{class\_create(name)}
\end{center}
Senin kodunda \texttt{class\_create("ahmet\_class")} -- \textbf{tek
argümanlı} -- kullanılıyor, yani sen kitaptan \textbf{iki nesil sonraki}
bir API'deyken (\texttt{owner} parametresi tamamen kaldırılmış).
İsimler değişse de \textbf{mekanizma birebir aynı}: sınıf oluştur
$\to$ \texttt{/sys/class/<isim>} dizini oluşsun $\to$ cihaz ekle
$\to$ \texttt{dev} dosyası ve hotplug olayı üretilsin $\to$ udev
\texttt{/dev} düğümünü yaratsın.
\end{ahmetkutu}

\subsection{Firmware Y\"ukleme (K\i sa Not)}

Donanımın çalışabilmesi için önce bir firmware dosyası yüklenmesi
gerekiyorsa, kitap firmware'i \textbf{kaynak koduna gömmeyi} açıkça
yanlış bulur (lisans sorunu + güncelleme zorluğu). Doğru yol:
\texttt{request\_firmware()} ile user space'ten çalışma zamanında
istemek -- bu çağrı \textbf{uyuyabilir} (process bağlamı gerektirir),
uyuyamayan bağlamlar için \texttt{request\_firmware\_nowait()}
(callback tabanlı) vardır.

\subsection*{B\"ol\"um 14 -- \"On Plana \c{C}\i kan Uyar\i lar}
\begin{itemize}[leftmargin=1.2cm]
  \item Bir kobject içeren struct'ı asla doğrudan \texttt{kfree} etme
  -- sadece \texttt{release} callback'i içinde, \texttt{kobject\_put}
  tetiklediğinde.
  \item sysfs dosyaları "her dosyada tek değer" kuralına uymalı.
  \item \texttt{show}/\texttt{store} fonksiyonları user space'ten gelen
  veriyi asla güvenilir kabul etmemeli -- geçersizse negatif hata dön.
  \item Kitaptaki \texttt{class\_simple\_*} isimlerini gördüğünde
  paniklemeyin -- bugünkü \texttt{class\_create}/\texttt{device\_create}'in
  doğrudan atalarıdır, aynı mekanizmayı farklı isimle sunar.
  \item Firmware'i asla kaynak koduna gömme -- hem lisans hem
  güncellenebilirlik sorunu yaratır.
\end{itemize}

\newpage

%=================================================================
\section*{Sonu\c{c}: Kitaptan ahmet.c'ye D\"onen Bir Yol Haritas\i}
\addcontentsline{toc}{section}{Sonu\c{c}: Kitaptan ahmet.c'ye D\"onen Bir Yol Haritas\i}
%=================================================================

Bu sekiz bölümü okuduktan sonra, kendi sürücünü aşamalı olarak
geliştirmek için önerilen sıra -- her adım tek bir bölümün tek bir
kavramına karşılık geliyor, yani her adımı bitirdiğinde \emph{hangi
bölümü uyguladığını} tam olarak bileceksin:

\begin{enumerate}[leftmargin=1.2cm]
  \item \textbf{(Bölüm 5)} \texttt{ahmet\_write}/\texttt{ahmet\_read}
  içine bir \texttt{struct mutex} ekle (modern karşılığı --
  \texttt{DECLARE\_MUTEX} değil). \texttt{ahmet\_open}'da değil,
  modülün \texttt{ahmet\_init}'inde, cihaz sisteme açılmadan
  \textbf{önce} ilklendir -- Bölüm 5'teki sıralama kuralını birebir
  uygula.
  \item \textbf{(Bölüm 3)} \texttt{ahmet\_write}'ı \texttt{*offset}'i
  kullanacak şekilde düzelt -- \texttt{scull\_write}'ın algoritmasını
  (kırpma + kopyalama + \texttt{*offset} güncelleme) birebir örnek al.
  \item \textbf{(Bölüm 6)} Basit bir \texttt{ioctl} komutu ekle
  (örneğin \texttt{AHMET\_IOCRESET} -- tamponu sıfırlayan, argümansız
  bir \texttt{\_IO} komutu). Bu, hem komut kodlamasını hem
  \texttt{.ioctl}/\texttt{.unlocked\_ioctl} metodunu eklemeyi
  gerektirecek.
  \item \textbf{(Bölüm 6)} \texttt{ahmet\_read}'i, veri yokken
  \texttt{wait\_event\_interruptible} ile bekleyecek şekilde
  değiştir -- \texttt{scull\_p\_read}'in "kilidi bırak, bekle, kilidi
  geri al, koşulu yeniden kontrol et" desenini uygula.
  \item \textbf{(Bölüm 2 + 8)} \texttt{BUFFER\_SIZE}'ı bir
  \texttt{module\_param}'a çevir, \texttt{kmalloc}/\texttt{kfree} ile
  dinamik ayır -- \texttt{ahmet\_init} process bağlamında çalıştığı
  için \texttt{GFP\_KERNEL} doğru bayraktır.
  \item \textbf{(Bölüm 14)} \texttt{class\_create}/\texttt{device\_create}
  çağrılarının kitaptaki \texttt{class\_simple\_create} ile aynı
  mekanizma olduğunu bilerek kodunu tekrar oku -- sysfs'teki
  \texttt{/sys/class/ahmet\_class/ahmet/dev} dosyasını
  (\texttt{cat /sys/class/ahmet\_class/ahmet/dev}) inceleyip udev'in
  bu dosyadan nasıl \texttt{/dev/ahmet}'i ürettiğini gözlemle.
  \item \textbf{(Bölüm 10, ileri seviye)} Staj kapsamında gerçek bir
  donanıma (GPIO, buton, sensör) bağlanırsan, \texttt{request\_irq}
  ile bir kesme işleyicisi ekle; handler'da sadece minimum işi yap ve
  Bölüm 6'daki \texttt{wait\_event\_interruptible} ile bloklanan bir
  \texttt{read}'i \texttt{wake\_up\_interruptible} ile uyandır --
  Bölüm 6 ve Bölüm 10'un birleştiği tam nokta budur.
  \item \textbf{(Bölüm 4)} Her adımdan sonra \texttt{dmesg} çıktısını
  ve gerekirse \texttt{strace}'i kullanarak davranışı doğrula --
  özellikle \texttt{ioctl} ve bloklayan \texttt{read} eklemelerinde,
  \texttt{strace} sana syscall seviyesinde tam olarak ne olduğunu
  gösterecek.
\end{enumerate}

\begin{notkutu}{Son not}
Bu belge kitabın yerini tutmaz -- burada anlatılan her kavramın
arkasında kitapta çok daha fazla ayrıntı, şekil ve alternatif örnek
var (özellikle \texttt{scullpipe}'ın tam kaynağı, \texttt{scullpriv}'in
liste yönetimi, \texttt{lddbus} örneği gibi). Buradaki amaç, kitabı
okurken hiçbir terimin ya da kod parçasının "havada kalmaması," ve
okuduğun her şeyi kendi \texttt{ahmet.c}'ne bir sonraki adım olarak
bağlayabilmendi. Yukarıdaki adımları sırayla uygularsan, staj sonunda
elinde sadece "çalışan" değil, \emph{LDD3'ün sekiz bölümünün her
birini bilinçli olarak uyguladığın} bir sürücü olacak.
\end{notkutu}

\end{document}
